\documentclass[10pt]{extarticle}
\usepackage[a4paper,margin=1.5cm]{geometry}
\usepackage[T1]{fontenc}
\usepackage{amsmath,amssymb}
\usepackage{graphicx}
\usepackage[export]{adjustbox}
\usepackage{multicol}
\usepackage{enumitem}
\usepackage{booktabs,array,tabularx}
\setlength{\columnsep}{0.7cm}
\setlength{\parindent}{0pt}
\setlength{\parskip}{3pt}
\pagestyle{empty}
\makeatletter
% figure registry (figdefs.tex is generated by the extract step)
\newcommand{\qfigdef}[4]{\@namedef{qf@file@#1}{#2}\@namedef{qf@w@#1}{#3}\@namedef{qf@wide@#1}{#4}}
\newcommand{\qf@check}[1]{\@ifundefined{qf@file@#1}{\PackageError{score2}{Unknown figure id #1}{}}{}}
\newcommand{\qf@img}[4][]{\includegraphics[width=#2\dimexpr\@nameuse{qf@w@#4}pt\relax,max width=#3,max height=0.45\textheight,#1]{\@nameuse{qf@file@#4}}}
% Full-width questions: finalize wraps every question that contains a wide figure in \qfullbegin...\qfullend,
% so the WHOLE question (stem, figures, options) is set across the page in reading order, then the two columns
% resume. The question's head - its number, stem and first figure (\qfullsplit marks the end of the head) -
% is kept on one page, so a number is never left at a page foot with its figure overleaf; the rest may break
% across pages normally, so no big gaps are left. Without a \qfullsplit, the whole question is kept together
% if it is under 45% of a page. Never write these yourself.
\newif\ifqf@full
\newif\ifqf@boxed
\newsavebox{\qf@box}
\newcommand{\qf@open}{\global\qf@fulltrue\par\global\qf@boxedtrue
  \setbox\qf@box\vbox\bgroup\hsize\textwidth\linewidth\textwidth}
\newcommand{\qf@emit}[1]{\par\egroup\global\qf@boxedfalse
  \ifdim\dimexpr\ht\qf@box+\dp\qf@box\relax<#1\textheight\noindent\usebox{\qf@box}\par\else\unvbox\qf@box\fi}
\newcommand{\qfullsplit}{\ifqf@boxed\qf@emit{0.85}\fi}
\newcommand{\qf@close}{\ifqf@boxed\qf@emit{0.45}\fi\global\qf@fullfalse}
\newcommand{\qfullbegin}{\end{multicols}\qf@open}
\newcommand{\qfullend}{\qf@close\begin{multicols}{2}}
\newcommand{\qfullstart}{\qf@open}
\newcommand{\qfulldone}{\qf@close}
% \qfig[scale]{id}: block figure, centred, directly where it is written (size set by the tool).
% Fallback if a wide figure is not inside a full-width question: break out of the columns in place - no floats.
\newcommand{\qfig}[2][1]{\qf@check{#2}%
  \ifqf@full\qf@blk{#1}{#2}\else
  \ifnum\@nameuse{qf@wide@#2}=1 \expandafter\@firstoftwo\else\expandafter\@secondoftwo\fi
  {\end{multicols}\begin{center}\qf@img{#1}{\textwidth}{#2}\end{center}\begin{multicols}{2}}%
  {\qf@blk{#1}{#2}}\fi}
\newcommand{\qf@blk}[2]{\par\smallskip{\centering\qf@img{#1}{\linewidth}{#2}\par}\smallskip}
% \qfigi[scale]{id}: inline figure (inside option lists / tables / minipages); never breaks columns.
\newcommand{\qfigi}[2][1]{\qf@check{#2}\qf@img[valign=c]{#1}{\linewidth}{#2}}
\makeatother
% structure macros - lint reads these to check completeness and figure placement
\newcommand{\papersec}[1]{\par\bigskip{\centering\Large\bfseries PAPER-#1\par}\medskip}
\newcommand{\partsec}[2]{\par\medskip\noindent{\large\bfseries PART-#1: #2}\par\smallskip}
\newcommand{\sect}[1]{\par\smallskip\noindent\textbf{#1}\par}
\newcommand{\qn}[1]{\par\smallskip\noindent\textbf{#1.}~}
\newcommand{\paraq}[2]{\par\smallskip\noindent\textbf{Paragraph for Questions #1 and #2}\par}
\newenvironment{qoptions}{\begin{itemize}[leftmargin=1.8em,itemsep=1pt,topsep=2pt,label={}]}{\end{itemize}}
\newcommand{\opt}[2]{\item[(#1)] #2}
% figure registry (generated by score2_tools.py finalize)
\qfigdef{p1_phy_s1_q02_1}{TNAS - Test 07 (20 Aug 2023)-img/p1_phy_s1_q02_1.png}{165.5}{0}
\qfigdef{p1_phy_s1_q05_1}{TNAS - Test 07 (20 Aug 2023)-img/p1_phy_s1_q05_1.png}{126.6}{0}
\qfigdef{p1_phy_s1_q07_1}{TNAS - Test 07 (20 Aug 2023)-img/p1_phy_s1_q07_1.png}{109.4}{0}
\qfigdef{p1_phy_s1_q11_1}{TNAS - Test 07 (20 Aug 2023)-img/p1_phy_s1_q11_1.png}{133.0}{0}
\qfigdef{p1_phy_s1_q12_1}{TNAS - Test 07 (20 Aug 2023)-img/p1_phy_s1_q12_1.png}{106.3}{0}
\qfigdef{p1_phy_s1_q15_1}{TNAS - Test 07 (20 Aug 2023)-img/p1_phy_s1_q15_1.png}{141.3}{0}
\qfigdef{p1_phy_s1_q17_1}{TNAS - Test 07 (20 Aug 2023)-img/p1_phy_s1_q17_1.png}{148.9}{0}
\qfigdef{p1_phy_s2_q04_1}{TNAS - Test 07 (20 Aug 2023)-img/p1_phy_s2_q04_1.png}{92.0}{0}
\qfigdef{p1_phy_s2_q08_1}{TNAS - Test 07 (20 Aug 2023)-img/p1_phy_s2_q08_1.png}{84.9}{0}
\qfigdef{p1_phy_s2_q09_1}{TNAS - Test 07 (20 Aug 2023)-img/p1_phy_s2_q09_1.png}{163.8}{0}
\qfigdef{p1_chem_s1_q03_1}{TNAS - Test 07 (20 Aug 2023)-img/p1_chem_s1_q03_1.png}{163.8}{0}
\qfigdef{p1_chem_s1_q15_1}{TNAS - Test 07 (20 Aug 2023)-img/p1_chem_s1_q15_1.png}{155.6}{0}
\qfigdef{p1_chem_s1_q15_2}{TNAS - Test 07 (20 Aug 2023)-img/p1_chem_s1_q15_2.png}{21.1}{0}
\qfigdef{p1_chem_s1_q15_3}{TNAS - Test 07 (20 Aug 2023)-img/p1_chem_s1_q15_3.png}{25.8}{0}
\qfigdef{p1_chem_s1_q15_4}{TNAS - Test 07 (20 Aug 2023)-img/p1_chem_s1_q15_4.png}{25.8}{0}
\qfigdef{p1_chem_s1_q15_5}{TNAS - Test 07 (20 Aug 2023)-img/p1_chem_s1_q15_5.png}{23.3}{0}
\qfigdef{p1_chem_s1_q17_1}{TNAS - Test 07 (20 Aug 2023)-img/p1_chem_s1_q17_1.png}{70.8}{0}
\qfigdef{p1_chem_s1_q17_2}{TNAS - Test 07 (20 Aug 2023)-img/p1_chem_s1_q17_2.png}{54.6}{0}
\qfigdef{p1_chem_s1_q17_3}{TNAS - Test 07 (20 Aug 2023)-img/p1_chem_s1_q17_3.png}{55.4}{0}
\qfigdef{p1_chem_s1_q17_4}{TNAS - Test 07 (20 Aug 2023)-img/p1_chem_s1_q17_4.png}{55.6}{0}
\qfigdef{p1_chem_s1_q18_A}{TNAS - Test 07 (20 Aug 2023)-img/p1_chem_s1_q18_A.png}{38.5}{0}
\qfigdef{p1_chem_s1_q18_C}{TNAS - Test 07 (20 Aug 2023)-img/p1_chem_s1_q18_C.png}{25.5}{0}
\qfigdef{p1_chem_s1_q18_D}{TNAS - Test 07 (20 Aug 2023)-img/p1_chem_s1_q18_D.png}{58.2}{0}
\qfigdef{p1_chem_s2_q04_1}{TNAS - Test 07 (20 Aug 2023)-img/p1_chem_s2_q04_1.png}{117.0}{0}
\qfigdef{p1_chem_s2_q06_1}{TNAS - Test 07 (20 Aug 2023)-img/p1_chem_s2_q06_1.png}{67.1}{0}
\qfigdef{p1_chem_s2_q07_1}{TNAS - Test 07 (20 Aug 2023)-img/p1_chem_s2_q07_1.png}{163.0}{0}
\qfigdef{p1_chem_s2_q07_2}{TNAS - Test 07 (20 Aug 2023)-img/p1_chem_s2_q07_2.png}{138.9}{0}
\qfigdef{p1_chem_s2_q08_1}{TNAS - Test 07 (20 Aug 2023)-img/p1_chem_s2_q08_1.png}{38.9}{0}
\qfigdef{p1_chem_s2_q09_1}{TNAS - Test 07 (20 Aug 2023)-img/p1_chem_s2_q09_1.png}{66.4}{0}
\qfigdef{p1_chem_s2_q09_2}{TNAS - Test 07 (20 Aug 2023)-img/p1_chem_s2_q09_2.png}{69.4}{0}
\qfigdef{p1_chem_s2_q10_1}{TNAS - Test 07 (20 Aug 2023)-img/p1_chem_s2_q10_1.png}{50.4}{0}
\qfigdef{p2_phy_s1_q01_1}{TNAS - Test 07 (20 Aug 2023)-img/p2_phy_s1_q01_1.png}{111.1}{0}
\qfigdef{p2_phy_s1_q02_1}{TNAS - Test 07 (20 Aug 2023)-img/p2_phy_s1_q02_1.png}{295.1}{1}
\qfigdef{p2_phy_s1_q03_1}{TNAS - Test 07 (20 Aug 2023)-img/p2_phy_s1_q03_1.png}{205.2}{0}
\qfigdef{p2_phy_s1_q05_1}{TNAS - Test 07 (20 Aug 2023)-img/p2_phy_s1_q05_1.png}{263.3}{1}
\qfigdef{p2_phy_s2_q07_1}{TNAS - Test 07 (20 Aug 2023)-img/p2_phy_s2_q07_1.png}{57.8}{0}
\qfigdef{p2_phy_s2_q07_2}{TNAS - Test 07 (20 Aug 2023)-img/p2_phy_s2_q07_2.png}{61.2}{0}
\qfigdef{p2_phy_s2_q07_3}{TNAS - Test 07 (20 Aug 2023)-img/p2_phy_s2_q07_3.png}{64.7}{0}
\qfigdef{p2_phy_s2_q07_4}{TNAS - Test 07 (20 Aug 2023)-img/p2_phy_s2_q07_4.png}{69.2}{0}
\qfigdef{p2_phy_s2_q07_5}{TNAS - Test 07 (20 Aug 2023)-img/p2_phy_s2_q07_5.png}{58.3}{0}
\qfigdef{p2_phy_s2_q08_1}{TNAS - Test 07 (20 Aug 2023)-img/p2_phy_s2_q08_1.png}{158.7}{0}
\qfigdef{p2_phy_s2_q09_1}{TNAS - Test 07 (20 Aug 2023)-img/p2_phy_s2_q09_1.png}{41.5}{0}
\qfigdef{p2_phy_s2_q09_2}{TNAS - Test 07 (20 Aug 2023)-img/p2_phy_s2_q09_2.png}{43.2}{0}
\qfigdef{p2_phy_s2_q09_3}{TNAS - Test 07 (20 Aug 2023)-img/p2_phy_s2_q09_3.png}{68.9}{0}
\qfigdef{p2_phy_s2_q09_4}{TNAS - Test 07 (20 Aug 2023)-img/p2_phy_s2_q09_4.png}{67.7}{0}
\qfigdef{p2_phy_s3_q02_1}{TNAS - Test 07 (20 Aug 2023)-img/p2_phy_s3_q02_1.png}{108.2}{0}
\qfigdef{p2_phy_s3_q03_1}{TNAS - Test 07 (20 Aug 2023)-img/p2_phy_s3_q03_1.png}{213.3}{0}
\qfigdef{p2_phy_s3_q04_1}{TNAS - Test 07 (20 Aug 2023)-img/p2_phy_s3_q04_1.png}{108.3}{0}
\qfigdef{p2_phy_s3_q06_1}{TNAS - Test 07 (20 Aug 2023)-img/p2_phy_s3_q06_1.png}{82.3}{0}
\qfigdef{p2_phy_s3_q07_1}{TNAS - Test 07 (20 Aug 2023)-img/p2_phy_s3_q07_1.png}{139.1}{0}
\qfigdef{p2_phy_s3_q08_1}{TNAS - Test 07 (20 Aug 2023)-img/p2_phy_s3_q08_1.png}{106.6}{0}
\qfigdef{p2_chem_s1_q05_1}{TNAS - Test 07 (20 Aug 2023)-img/p2_chem_s1_q05_1.png}{265.4}{1}
\qfigdef{p2_chem_s1_q06_1}{TNAS - Test 07 (20 Aug 2023)-img/p2_chem_s1_q06_1.png}{168.0}{0}
\qfigdef{p2_chem_s1_q06_2}{TNAS - Test 07 (20 Aug 2023)-img/p2_chem_s1_q06_2.png}{159.0}{0}
\qfigdef{p2_chem_s2_q09_1}{TNAS - Test 07 (20 Aug 2023)-img/p2_chem_s2_q09_1.png}{57.0}{0}
\qfigdef{p2_chem_s2_q09_2}{TNAS - Test 07 (20 Aug 2023)-img/p2_chem_s2_q09_2.png}{57.1}{0}
\qfigdef{p2_chem_s3_q03_1}{TNAS - Test 07 (20 Aug 2023)-img/p2_chem_s3_q03_1.png}{108.3}{0}
\qfigdef{p2_chem_s3_q04_1}{TNAS - Test 07 (20 Aug 2023)-img/p2_chem_s3_q04_1.png}{132.2}{0}
\qfigdef{p2_chem_s3_q06_1}{TNAS - Test 07 (20 Aug 2023)-img/p2_chem_s3_q06_1.png}{296.8}{1}
\qfigdef{p2_chem_s3_q07_1}{TNAS - Test 07 (20 Aug 2023)-img/p2_chem_s3_q07_1.png}{245.2}{0}
\qfigdef{p2_chem_s3_q07_2}{TNAS - Test 07 (20 Aug 2023)-img/p2_chem_s3_q07_2.png}{235.8}{0}
\qfigdef{p2_chem_s3_q07_3}{TNAS - Test 07 (20 Aug 2023)-img/p2_chem_s3_q07_3.png}{190.1}{0}
\qfigdef{p2_chem_s3_q08_1}{TNAS - Test 07 (20 Aug 2023)-img/p2_chem_s3_q08_1.png}{56.6}{0}
\raggedcolumns
\makeatletter\@beginparpenalty=10000\makeatother
\newcolumntype{L}[1]{>{\raggedright\arraybackslash}p{#1}}
\begin{document}
\begin{center}{\Large\bfseries TNAS --- Test 07 (20 Aug 2023)}\end{center}
\begin{multicols}{2}
\raggedcolumns
\papersec{1}
\partsec{1}{PHYSICS}
\sect{SECTION-I}
\qn{1} A hypothetical e-nucleus system obeys bohr's postulates and it's potential energy is expressed as $U=-\dfrac{K}{4r^4}$, where $K$ is positive constant then dependency of kinetic energy on principal quantum number $n$ is:
\begin{qoptions}
\opt{A}{$KE\times1/n^2$}
\opt{B}{$KE\times n^2$}
\opt{C}{$KE\times n^{-4}$}
\opt{D}{None}
\end{qoptions}
\qn{2} For x-rays spectrum is as given for three different atoms for which atomic number are $z_1$, $z_2$ \& $z_3$, while the accelerating potentials are $v_1$, $v_2$ \& $v_3$ respectively then, choose the correct relation.
\qfig{p1_phy_s1_q02_1}
(1) $z_1>z_2>z_3$ \quad (2) $v_1>v_2>v_3$ \quad (3) $z_3>z_1>z_2$ \quad (4) $v_3>v_1>v_2$
\begin{qoptions}
\opt{A}{1 \& 2}
\opt{B}{2 \& 3}
\opt{C}{3 \& 4}
\opt{D}{None}
\end{qoptions}
\qn{3} \textbf{Assertion :} When hydrogen atom is raised from the ground state to an excited state the kinetic energy increases.\par
and\par
\textbf{Reason :} The velocity of electron in a quantum state is inversely proportional to the quantum number of the state.
\begin{qoptions}
\opt{A}{Assertion is true, Reason is true and Reason is correct explanation for Assertion.}
\opt{B}{Assertion is true, Reason is true and Reason is NOT correct explanation for Assertion.}
\opt{C}{Assertion is true, Reason is false.}
\opt{D}{Assertion is false, Reason is true.}
\end{qoptions}
\qn{4} A particle of mass $m$ is moving in potential energy $U=\dfrac12kx^2+\dfrac{\alpha}{2x^2}$ where $k$ and $\alpha$ are positive constants. The angular frequency of small oscillations for the simple harmonic motion of the particle about a stable equilibrium position is
\begin{qoptions}
\opt{A}{$\sqrt{\dfrac km}$}
\opt{B}{$\sqrt{\dfrac{2k}m}$}
\opt{C}{$2\sqrt{\dfrac km}$}
\opt{D}{$4\sqrt{\dfrac km}$}
\end{qoptions}
\qn{5} A symmetrical block of mass $M$ having a notch of hemispherical shape of radius $R$ is placed on a horizontal surface. A small particle of mass $m$ is placed inside the hemispherical notch and can slide without friction. Find the time period of small oscillation of the particle and the block. The acceleration due to gravity is $g$.
\qfig{p1_phy_s1_q05_1}
\begin{qoptions}
\opt{A}{$2\pi\sqrt{\dfrac{R}{g\left[1+\frac mM\right]}}$}
\opt{B}{$2\pi\sqrt{\dfrac{R}{g\left[1+\frac Mm\right]}}$}
\opt{C}{$2\pi\sqrt{\dfrac Rg\left[1+\dfrac mM\right]}$}
\opt{D}{$2\pi\sqrt{\dfrac Rg\left[1+\dfrac Mm\right]}$}
\end{qoptions}
\qn{6} When a telescope is adjusted for parallel light, the distance of the objective from the eyepiece is 100 cm for normal adjustment. The magnifying power of the telescope in this case is 9. If an old man cannot see beyond 90 cm and wishes to use the telescope then he will have to reduce the tube length by
\begin{qoptions}
\opt{A}{1.5 cm}
\opt{B}{0.5 cm}
\opt{C}{1 cm}
\opt{D}{2 cm}
\end{qoptions}
\qn{7} A light ray is incident on a prism of prism angle $60^\circ$ and refractive index $\sqrt3$ as shown in figure. The angle of emergence $e$ is given by :
\qfig{p1_phy_s1_q07_1}
\begin{qoptions}
\opt{A}{$30^\circ$}
\opt{B}{$45^\circ$}
\opt{C}{$90^\circ$}
\opt{D}{$60^\circ$}
\end{qoptions}
\qn{8} Chromatic aberration of a lens can be corrected by :
\begin{qoptions}
\opt{A}{providing different suitable curvature of its two surfaces}
\opt{B}{proper polishing of its two surface}
\opt{C}{suitably combining it with another lens}
\opt{D}{None of these}
\end{qoptions}
\qn{9} For a concave mirror. The correct matching of items on the left with those on the right is\par
{\footnotesize\setlength{\tabcolsep}{2pt}\begin{tabular}{|c|p{0.3\linewidth}|c|p{0.2\linewidth}|}
\hline
\multicolumn{2}{|c|}{\textbf{List-I}} & \multicolumn{2}{c|}{\textbf{List-II}}\\ \hline
(P) & Between centre of curvature and focus & (1) & enlarged, virtual, erect\\ \hline
(Q) & Beyond the centre of curvature (away from the mirror) & (2) & enlarged, real, inverted\\ \hline
(R) & Between focus and pole & (3) & same size, real, inverted\\ \hline
(S) & At the centre of curvature & (4) & diminished, real, inverted\\ \hline
\end{tabular}}\par
\begin{qoptions}
\opt{A}{(P)-(2); (Q)-(3) ; (R)-(1) ; (S)-(4)}
\opt{B}{(P)-(3); (Q)-(2) ; (R)-(1) ; (S)-(4)}
\opt{C}{(P)-(2) ; (Q)-(1) ; (R)-(3) ; (S)-(4)}
\opt{D}{(P)-(2) ; (Q)-(4) ; (R)-(1) ; (S)-(3)}
\end{qoptions}
\qn{10} A far sighted person has a near point 60 cm away. What should be the power of a lens he should use for eye glasses so that he can read a book at a distance of 25 cm?
\begin{qoptions}
\opt{A}{1.33D}
\opt{B}{2.33D}
\opt{C}{$-1.33$D}
\opt{D}{$-2.33$D}
\end{qoptions}
\qn{11} As shown in figure, a thin convex lens of focal length 20 cm is inclined at an angle of $53^\circ$ with $x$-axis. A point object is placed on the $x$-axis at a distance of 100 cm from the centre of the lens. Taking the origin as the centre of lens, find the $x$-coordinate of the image.
\qfig{p1_phy_s1_q11_1}
\begin{qoptions}
\opt{A}{26.7 cm}
\opt{B}{25 cm}
\opt{C}{33.3 cm}
\opt{D}{20 cm}
\end{qoptions}
\qn{12} Two rays are incident on a spherical concave mirror of radius $R=5$ cm and rays are parallel to a line $AB$ which passes through the centre $C$. The rays are at perpendicular distances 3 cm and 4 cm from $AB$ respectively as shown. Find the distance between the points at which these rays intersect the line $AB$ after reflection :
\qfig{p1_phy_s1_q12_1}
\begin{qoptions}
\opt{A}{$\dfrac{25}{24}$ cm}
\opt{B}{$\dfrac{24}{25}$ cm}
\opt{C}{zero}
\opt{D}{$\dfrac13$ cm}
\end{qoptions}
\qn{13} A xenon arc lamp is covered with an interference filter that only transmits light of 400 nm wavelength. The transmitted light strikes a metal surface, a stream of electrons emerges from the metal. The filter is then replaced with one transmitting at 300 nm and the lamp adjusted so that the intensity of striking the surface is the same as it was for the 400 nm light. With the 300 nm light :-
\begin{qoptions}
\opt{A}{More electrons are emitted in a given time interval}
\opt{B}{The electrons which are emitted are more energetic}
\opt{C}{Both are true}
\opt{D}{Both are false}
\end{qoptions}
\qn{14} Which of the following statements is \textbf{INCORRECT} ?
\begin{qoptions}
\opt{A}{The debroglie wavelength of electrons in Bohr orbit is quantized.}
\opt{B}{Electron de-broglie wavelength depends on the kinetic energy of electron.}
\opt{C}{Energy of electrons is quantised in $\mathrm{Li^{++}}$ ion.}
\opt{D}{Electrons can not exist inside nucleus because they are attracted by protons inside nucleus.}
\end{qoptions}
\qn{15} The quantum levels of atoms of hydrogen and deuterium are only approximately the same\par
(the difference between the two systems of levels is exaggerated in the figure).
\qfig{p1_phy_s1_q15_1}
\textbf{Statement-1 :} A corresponds to quantum levels of hydrogen atom and B to quantum levels of deuterium atom.\par
\textbf{Statement-2 :} The mass of H nuclei is lesser than the mass of deuterium nuclei.
\begin{qoptions}
\opt{A}{Statement-1 is true, statement-2 is true and statement-2 is correct explanation for statement-1.}
\opt{B}{Statement-1 is true, statement-2 is true and statement-2 is NOT the correct explanation for statement$-1$.}
\opt{C}{Statement-1 is true, statement-2 is false.}
\opt{D}{Statement-1 is false, statement-2 is true.}
\end{qoptions}
\qn{16} In a cricket match, batsman hit the ball (mass $=m$) high in the air and fielding side player catch the ball near boundary. De broglie wavelength of ball just after hit from bat is $\lambda_1$ and just before the catch is $\lambda_2$. Assume height at which ball is hit by bat and height at which ball is catch are same. Find total work done by air-resistance (Mechanical energy is loss only due to air-resistance)
\begin{qoptions}
\opt{A}{$\dfrac{h^2}m\left(\dfrac1{\lambda_1}-\dfrac1{\lambda_2}\right)$}
\opt{B}{$\dfrac{h^2}{2m}\left(\dfrac1{\lambda_1^2}-\dfrac1{\lambda_2^2}\right)$}
\opt{C}{$\dfrac{h^2}{2m}\left(\dfrac{\lambda_1\lambda_2}{\lambda_1-\lambda_2}\right)$}
\opt{D}{$\dfrac{h^2}{2m}\left(\dfrac1{\lambda_1}-\dfrac1{-\lambda_2}\right)$}
\end{qoptions}
\qn{17} From last few decades Scientist Search alternative option to move in space. One idea is to build a boat (mass $=m$) with large canvas (area $A$) and accelerate boat by radiation force exerted by solar radiation. Find acceleration of boat when it is at distance $r$ from sun (power of sun $=P$ ; assume rays are normally incident on canvas and canvas is 100\% reflector of light)
\qfig{p1_phy_s1_q17_1}
\begin{qoptions}
\opt{A}{$\dfrac{PA}{4\pi r^2cm}$}
\opt{B}{$\dfrac{PAc}{4\pi r^2m}$}
\opt{C}{$\dfrac{PA}{2\pi r^2cm}$}
\opt{D}{$\dfrac{PAc}{2\pi r^2m}$}
\end{qoptions}
\qn{18} Both the frequency and the intensity of a beam of light falling on the surface of photoelectric material are increased by a factor of two. This will:
\begin{qoptions}
\opt{A}{increase both, the maximum kinetic energy of the photo-electrons, as well as photoelectric saturation current by a factor of two.}
\opt{B}{increase the maximum kinetic energy of the photo-electrons by a factor greater than two and would increase the photoelectric saturation current by a factor of two.}
\opt{C}{increase the maximum kinetic energy of the photoelectrons by a factor greater than two but and will have no effect on the magnitude of the photoelectric saturation current produced.}
\opt{D}{increase the maximum kinetic energy of the emitted photo-electrons by a factor of two but will have no effect on saturation photoelectric current.}
\end{qoptions}
\qn{19} Internal conversion is the process whereby an excited nucleus transfers its energy directly to one of the most tightly bound atomic electrons, causing the electron to be ejected from the atom and leaving the atom in an excited state. The most probable process after an internal conversion from an atom with a high atomic number is that the
\begin{qoptions}
\opt{A}{atom returns to ground state through inelastic collisions with other atoms}
\opt{B}{atom emits one or several X-rays}
\opt{C}{nucleus emits a $\gamma$-ray}
\opt{D}{nucleus emits an electron}
\end{qoptions}
\qn{20} In a hypothetical atom, a negatively charged particle having a charge of magnitude $3e$ and mass $3m$ revolves around a proton. Here, $e$ is the electronic charge and $m$ is the electronic mass. Mass of proton may be assumed to be much larger than that of the negatively charged particle, thus the proton is at rest. This ``atom'' obeys Bohr's postulate of quantization of angular momentum, that is $mvr=n\left(\dfrac h{2\pi}\right)$. It is given that for the first Bohr orbit of hydrogen atom: radius of orbit is $r_0$, speed of electron is $V_0$, and total energy is $-E_0$. The momentum of an emitted photon when it makes a transition from the second excited state to ground state, is
\begin{qoptions}
\opt{A}{$\dfrac{8E_0}c$}
\opt{B}{$\dfrac{E_0}c$}
\opt{C}{$\dfrac{24E_0}c$}
\opt{D}{None of these}
\end{qoptions}
\sect{SECTION-II}
\qn{1} In an electromagnetic wave, the amplitude of electric field is 10 V/m. The frequency of wave is $5\times10^{14}$ Hz, the wave is propagating along $z$-axis, then total average energy density of E.M. wave is $\alpha\times10^{-10}$ J/m$^3$, find the value of $\alpha$ to the nearest integer.
\qn{2} Assume that a plane wave in air with an electric field $\vec E=10\cos\left(\omega t-3x-\sqrt3z\right)\hat a_y$ V/m is incident on a non-magnetic dielectric slab of relative permittivity 3 which covers the region $z>0$. The angle of transmission in the dielectric slab Is \underline{\hspace{1cm}} degrees
\qn{3} Converging lens of diameter of aperture 5 cm with a focal length $F=50$ cm was cut in half and spaced at a distance $d=5$ mm. A point source of light $S$ is located at a distance $a=75$ cm from the lens symmetrically. At what distance $\ell$ (in cm) from the lens will interference begin ? The gap between the halves of the lens is closed by an opaque body.
\qn{4} A very fine hair-like fiber is stuck between two microscope slides. As a result of this, there is a wedge of air between them. When the slides are illuminated normally (from above) with light of wavelength 550 nm, bright and dark interference bands are formed. The fiber is seen to lie at the position of fifth dark band counting from the common edge. What is the diameter (in nm) of the fiber?
\qfig{p1_phy_s2_q04_1}
\qn{5} An unpolarized light beam is incident on the polarizer of a polarization experiment and the intensity of light beam emerging from the analyzer is measured as 100 Lumens. Now, if the analyzer is rotated around the horizontal axis (direction of light) by $30^\circ$ in clockwise direction, the intensity of emerging light will be \underline{\hspace{1cm}}
\qn{6} The $K_\alpha$ X-ray of molybdenum has wavelength 0.071 nm. If the energy of a molybdenum atoms with a K electron knocked out is 27.5 keV, the energy of this atom when an L electron is knocked out will be \underline{\hspace{1cm}} keV. (Round off to the nearest integer) [$h=4.14\times10^{-15}$ eVs, $c=3\times10^8$ ms$^{-1}$]
\qn{7} A damped harmonic oscillator has a frequency of 5 oscillations per second. The amplitude drops to half its value for every 10 oscillations. The time (in sec) it will take to $\dfrac1{1000}$ of the original amplitude is close to - ($\ell n\,10=2.3$)
\qn{8} A rectangular conducting loop of length 4 cm and width 2 cm is in the $xy$-plane, as shown in the figure. It is being moved away from a thin and long conducting wire along the direction $\dfrac{\sqrt3}2\hat x+\dfrac12\hat y$ with a constant speed $v$. The wire is carrying a steady current $I=10$ A in the positive $x$-direction.\par
A current of 10 $\mu$A flows through the loop when it is at a distance $d=4$ cm from the wire. If the resistance of the loop is 0.1 $\Omega$, then the value of $v$ is \underline{\hspace{1cm}} ms$^{-1}$. [Given: The permeability of free space $\mu_0=4\pi\times10^{-7}$ NA$^{-2}$]
\qfig{p1_phy_s2_q08_1}
\qn{9} A wire is made by welding together two metals having different densities. Figure shows a 2.00m long section of wire centered on the junction, but the wire extends much farther in both directions. The wire is placed under 2250 N tension, then a 1500 Hz wave with an amplitude of 3.00 mm is sent down the wire. How many wavelengths (complete cycles) of the wave are in this 2.00 m long section of the wire?
\qfig{p1_phy_s2_q09_1}
\qn{10} A plane electromagnetic wave of frequency $5\times10^{14}$ Hz and amplitude $10^3$ V/m travelling in a homogeneous dielectric medium of dielectric constant 1.69 is incident normally at the interface with a second dielectric medium of dielectric constant 2.25 . The ratio of the amplitude of the transmitted wave to that of the incident wave is :- (to the nearest integer)
\partsec{2}{CHEMISTRY}
\sect{SECTION-I}
\qn{1} Select the property which have higher value for $\mathrm{H_2O}$ when compared to $\mathrm{D_2O}$.
\begin{qoptions}
\opt{A}{Density}
\opt{B}{Boiling point}
\opt{C}{Viscosity}
\opt{D}{Dielectric constant}
\end{qoptions}
\qn{2} Select the correct statement regarding hydrogen as a fuel.
\begin{qoptions}
\opt{A}{The energy released on combustion per gram for hydrogen is less than LPG but greater than petrol.}
\opt{B}{The pollutants that will be released will be oxides of nitrogen for which there is no way of reducing it.}
\opt{C}{The cylinder of compressed hydrogen would weigh much lighter than the tank of petrol as hydrogen is lightest gas.}
\opt{D}{Expensive insulted tanks made from metal alloys like $\mathrm{NaNi_5}$ , $\mathrm{Ti{-}TiH_2}$ , $\mathrm{Mg{-}MgH_2}$ etc are in use for storage of dihydrogen in small quantities.}
\end{qoptions}
\qn{3}
\qfig{p1_chem_s1_q03_1}
Select the \textbf{INCORRECT} option :-
\begin{qoptions}
\opt{A}{$C$ can be obtained from $A$ by treating with Zn in the presence of HCl}
\opt{B}{$C$ can be obtained from $B$ by using $\mathrm{H_2O_2}$ solution}
\opt{C}{$C$ can be obtained from $A$ by fusing with $\mathrm{Na_2CO_3}$}
\opt{D}{$C$ can be dissolved in aqua regia}
\end{qoptions}
\qn{4} When the volatile complex formed in the Mond's process decomposes, the bond length of C$-$O.
\begin{qoptions}
\opt{A}{Increases}
\opt{B}{Decreases}
\opt{C}{Remains same}
\opt{D}{Can't be predicted}
\end{qoptions}
\qn{5} Which of the following alloys do not contain copper?
\begin{qoptions}
\opt{A}{German silver}
\opt{B}{Brass}
\opt{C}{Bronze}
\opt{D}{Solder}
\end{qoptions}
\qn{6} $\mathrm{NiCl_2[P(C_2H_5)_2(C_6H_5)]_2}$ exhibits temperature dependent magnetic behaviour (paramagnetic/diamagnetic). The difference in magnetic moment in both the state is :-
\begin{qoptions}
\opt{A}{1.73 B.M.}
\opt{B}{2.83 B.M.}
\opt{C}{3.87 B.M.}
\opt{D}{0 B.M.}
\end{qoptions}
\qn{7} The boiling point of an azeotropic mixture of phenol-aniline is higher than that of both phenol and aniline, then correct statement is :
\begin{qoptions}
\opt{A}{The mixture shows positive deviation from Raoult's law.}
\opt{B}{Change in enthalpy during mixing these liquids is less than zero.}
\opt{C}{Change in entropy during mixing these liquids is less than zero.}
\opt{D}{The mixture follows Raoult's law.}
\end{qoptions}
\qn{8} During a reversible adiabatic process, the pressure of a gas is found to be proportional to cube of its absolute temperature. The heat capacity ratio ($C_P/C_V$) for the gas is :
\begin{qoptions}
\opt{A}{3/2}
\opt{B}{7/2}
\opt{C}{5/3}
\opt{D}{9/7}
\end{qoptions}
\qn{9} In an exothermic reaction $A(g)\rightleftharpoons B(g)$, the activation energy of reverse reaction is twice that of forward reaction. If enthalpy of reaction is $-90$ kJ/mol, the activation energy of the reverse reaction is :
\begin{qoptions}
\opt{A}{60 kJ/mol}
\opt{B}{45 kJ/mol}
\opt{C}{180 kJ/mol}
\opt{D}{90 kJ/mol}
\end{qoptions}
\qn{10} For 1 mole of an ideal monoatomic gas on moving from one state to another state, the temperature of gas is doubled and its pressure become 4 times. Then entropy change in the process will be :
\begin{qoptions}
\opt{A}{$\dfrac{R\ln2}{2}$}
\opt{B}{$\dfrac{9R\ln2}{2}$}
\opt{C}{$2R\ln2$}
\opt{D}{$\dfrac{5R\ln2}{2}$}
\end{qoptions}
\qn{11} Which one of the following statements is incorrect for NaCl crystal :\par
[$a$ = Edge length of unit cell]
\begin{qoptions}
\opt{A}{Closest distance between cation-cation is $\dfrac a{\sqrt2}$.}
\opt{B}{$\mathrm{Na^+}$ ions are present at body centre and edge centres of FCC unit cell made of $\mathrm{Cl^-}$ ions.}
\opt{C}{Each $\mathrm{Na^+}$ ion is closely surrounded by 6 other $\mathrm{Na^+}$ ions.}
\opt{D}{Closest distance between cation-anion is $\dfrac a2$.}
\end{qoptions}
\qn{12} \textbf{Assertion :} Entropy of an ideal gas increases during isothermal free expansion of the gas.\par
\textbf{Reason :} Isothermal free expansion of an ideal gas is also adiabatic free expansion.
\begin{qoptions}
\opt{A}{If both assertion and reason are true and the reason is a correct explanation of assertion.}
\opt{B}{If both assertion and reason are true but reason is not a correct explanation of assertion.}
\opt{C}{If assertion is true but the reason is false.}
\opt{D}{If assertion is false but the reason is true.}
\end{qoptions}
\qn{13} Which of the following equation must be used for the exact calculation of $[\mathrm{H^+}]$ of an aqueous HCl solution at any concentration.
\begin{qoptions}
\opt{A}{$[\mathrm{H^+}]=[\mathrm{HCl}]_0+K_W$}
\opt{B}{$[\mathrm{H^+}]=[\mathrm{HCl}]_0+\dfrac{K_w}{[\mathrm{H^+}]}$}
\opt{C}{$[\mathrm{H^+}]=[\mathrm{HCl}]_0+\dfrac{K_w}{[\mathrm{OH^-}]}$}
\opt{D}{$[\mathrm{H^+}]=[\mathrm{HCl}]_0-\dfrac{K_w}{[\mathrm{OH^-}]}$}
\end{qoptions}
\qn{14}
{\scriptsize\setlength{\tabcolsep}{2pt}\begin{tabular}{|c|p{0.2\linewidth}|c|p{0.34\linewidth}|}
\hline
\multicolumn{2}{|c|}{\textbf{List-I} Aq. metal ion} & \multicolumn{2}{c|}{\textbf{List-II} Observation with given reagent}\\ \hline
(P) & $\mathrm{Fe^{2+}}$ ion & (1) & blue ppt. with $\mathrm{K_3[Fe(CN)_6]}$\\ \hline
(Q) & $\mathrm{Fe^{3+}}$ ion & (2) & orange ppt. with $\mathrm{K_3[Fe(CN)_6]}$\\ \hline
(R) & $\mathrm{Cu^{2+}}$ ion & (3) & green ppt. with $\mathrm{K_3[Fe(CN)_6]}$\\ \hline
(S) & $\mathrm{Ag^{+}}$ ion & (4) & brown colour with $\mathrm{K_3[Fe(CN)_6]}$\\ \hline
\end{tabular}}\par
\begin{qoptions}
\opt{A}{$P\to1,Q\to4;R\to3;S\to2$}
\opt{B}{$P\to2,Q\to3;R\to1;S\to4$}
\opt{C}{$P\to1,Q\to3;R\to2;S\to4$}
\opt{D}{$P\to4,Q\to1;R\to3;S\to2$}
\end{qoptions}
\qn{15}
\qfig{p1_chem_s1_q15_1}
Final product is :
\begin{qoptions}
\opt{A}{\qfigi{p1_chem_s1_q15_2}}
\opt{B}{\qfigi{p1_chem_s1_q15_3}}
\opt{C}{\qfigi{p1_chem_s1_q15_4}}
\opt{D}{\qfigi{p1_chem_s1_q15_5}}
\end{qoptions}
\qn{16} Reaction of acidified permanganate solution with $\mathrm{(CH_3)_2C{=}CH{-}CH_2{-}COOH}$ gives (A) and (B) compounds. A gives postive Iodoform test\par
Choose incorrect statement
\begin{qoptions}
\opt{A}{Compound (A) give $2^\circ$--alcohol by $\mathrm{LiAlH_4}$}
\opt{B}{Compound (B) is heated to give acetic acid}
\opt{C}{Compound (A) gives $\alpha,\beta$ unsaturated carbonyl compound with $\mathrm{Ba(OH)_2}/\Delta$.}
\opt{D}{Compound (B) is stronger acid than oxalic acid}
\end{qoptions}
\qn{17} Which of the following statement is incorrect?
\begin{qoptions}
\opt{A}{With $\mathrm{SeO_2}$, $\mathrm{CH_3CHO}$ gives glyoxal where acetone gives methylglyoxal}
\opt{B}{Acetone give yellow ppt of Iodoform on treatment $\mathrm{I_2}$/NaOH}
\opt{C}{Conc. KOH changes \qfigi{p1_chem_s1_q17_1}\qfigi{p1_chem_s1_q17_2}}
\opt{D}{$\mathrm{(CH_3)_2\,CuLi}$ changes \qfigi{p1_chem_s1_q17_3}\qfigi{p1_chem_s1_q17_4}}
\end{qoptions}
\qn{18} Which of the following reaction can not form major product as propane
\begin{qoptions}
\opt{A}{\qfigi{p1_chem_s1_q18_A}$\xrightarrow[\Delta]{\ominus\mathrm{OH}}$}
\opt{B}{$\mathrm{CH_3{-}CH{=}CH_2}\xrightarrow[\mathrm{H_2O_2}]{\mathrm{N_2H_4}}$}
\opt{C}{\qfigi{p1_chem_s1_q18_C}$\xrightarrow[\Delta]{\mathrm{H_2/Raney\ Ni}}$}
\opt{D}{\qfigi{p1_chem_s1_q18_D}}
\end{qoptions}
\qn{19} Consider following reaction sequence\par
$\mathrm{CH_3{-}N{=}C{=}O}\xrightarrow{\mathrm{H_3O^\oplus}}P+Q$\par
$(Q)\xrightarrow[\mathrm{H^\oplus}]{\mathrm{EtMgBr}}(S)\xrightarrow[\text{(ii) Aq.KOH}]{\text{(i) Red P + }\mathrm{Br_2}}(T)\ \text{(Major)}\xrightarrow{\Delta}(U)$\par
If aqueous solution of $Q$ is more acidic than $P$, then choose incorrect option:
\begin{qoptions}
\opt{A}{$Q$ turns lime water milky}
\opt{B}{$P$ gives foul smell with $\mathrm{CHCl_3}$ + KOH / $\Delta$}
\opt{C}{Degree of unsaturation of `$U$' is two.}
\opt{D}{$U$ contains total three stereoisomer}
\end{qoptions}
\qn{20} \textbf{Statement-1 :} Carboxylic acid on Reduction with $\mathrm{B_2H_6}$ / THF gives alcohol as product.\par
\textbf{Statement-2 :} Less the electron density on carboxylic acid derivative more will it's reactivity towards $\mathrm{B_2H_6}$ / THF.\par
Choose correct statement.
\begin{qoptions}
\opt{A}{Statement-1 is true, statement-2 is true and statement-2 is correct explanation for statement-1.}
\opt{B}{Statement-1 is true, statement-2 is true and statement-2 is NOT the correct explanation for statement-1.}
\opt{C}{Statement-1 is true, statement-2 is false.}
\opt{D}{Statement-1 is false, statement-2 is true.}
\end{qoptions}
\sect{SECTION-II}
\qn{1} When one mole $\mathrm{[Ni(H_2O)_6]^{2+}}$ complex ($R$) reacts with excess of ethylene diammine, a complex ($A$) is formed.\par
Total number of stereoisomers in complex ($A$) is $=a$.\par
Total number of unpaired electron present in complex ($A$) is $=b$\par
Change in crystal field stabilisation energy in both the complexes $=c$\par
Calculate $(a\times b)+c$
\qn{2} Number of cations from the following showing different colour in borax bead test in oxidising and reducing flame.\par
$\mathrm{Cu^{2+}}$ , $\mathrm{Fe^{3+}}$ , $\mathrm{Cr^{3+}}$ , $\mathrm{Ni^{2+}}$ , $\mathrm{Mn^{2+}}$
\qn{3} How many of the following metals are present as anode mud during the electrolytic refining of copper.\par
Antimony Selenium , Tellurium , Silver, Gold, Platinum
\qn{4} One mole of an ideal monoatomic gas undergo the following cyclic process.
\qfig{p1_chem_s2_q04_1}
If magnitude of work (in cal) in cyclic process [ABCDA] is `$W$' cal. Then find value of $\left(\dfrac W{70}\right)$ ? (Use : $\ln2=0.7$)
\qn{5} $\Delta_CH^\circ$ of graphite $=-95$ kcal/mol\par
$\Delta_CH^\circ$ of $\mathrm{H_2(g)}=-68$ kcal/mol\par
$\Delta_fH^\circ$ of $\mathrm{C_4H_{10}(g)}=-42.9$ kcal/mol\par
Calculate magnitude of standard heat of combustion of butane gas at constant temperature of $27^\circ$C and constant volume. ($R=2$ cal/mol-K)
\qn{6} In a container of volume 2 dm$^3$, $A(g)$ is allowed to established equilibrium as $A(g)\rightleftharpoons nB(g)$. At equilibrium 5 mole $A$ and 3 mole $B$ are present. Variation of degree of dissociation ($\alpha$) with vapour density is shown by following graph.
\qfig{p1_chem_s2_q06_1}
where, $D_T$ = Theoretical vapour density, $D_O$ = Observed vapour density.\par
Calculate ($\mathbf{K_c\times10}$) for reaction, $A(g)\rightleftharpoons nB(g)$.
\qn{7} Find number of reaction giving correct major product?
\qfig{p1_chem_s2_q07_1}
(5) $\mathrm{CH_3{-}CH_2{-}MgBr}\xrightarrow[\text{(2) }\mathrm{H_3O^+}]{\text{(1) }\mathrm{CO_2}}\mathrm{CH_3{-}CH_2{-}COOH}$\par
\qfig{p1_chem_s2_q07_2}
\qn{8} The number of aldol reaction(s) that occurs in the given transformation is\par
$\mathrm{CH_3CHO}+4\mathrm{HCHO}\xrightarrow{\text{Conc. aq. NaOH}}$ \qfigi{p1_chem_s2_q08_1}
\qn{9} \qfigi{p1_chem_s2_q09_1}\ $X$\ \qfigi{p1_chem_s2_q09_2}\par
$Y$(Aromatic hydrocarbon)\par
In `$Y$' total number of sigma bond formed by overlapping of ($sp^2-sp^2$) orbital(s) $=$ ($P$)\par
Total number of $\pi e^-$ in `$Y$' $=$ ($Q$)\par
value of $|P-Q|$ is
\qn{10} Consider following reaction\par
$\mathrm{H_3C{-}CH_2{-}Cl}\xrightarrow[\text{(ii)}\mathrm{CH_3Li}]{\text{(i)}\mathrm{P(Ph)_3}}(A)$ $+$ \qfigi{p1_chem_s2_q10_1}\par
$(B)\xrightarrow{\mathrm{O_3/Zn}}(C+D)\xrightarrow[5^\circ\mathrm{C}]{\text{dil.NaOH}}$\par
Find total number of final aldol product(s) of given reaction sequence including stereo.
\partsec{3}{MATHEMATICS}
\sect{SECTION-I}
\qn{1} Let $A(-1,0)$ and $B(2,0)$ be two points on the $x$-axis. A point `$M$' is moving in $xy$-plane (other than $x$-axis) in such a way that $\angle MBA=2\angle MAB$, then the point `$M$' moves along a conic whose length of latus-rectum equals
\begin{qoptions}
\opt{A}{4}
\opt{B}{6}
\opt{C}{$\dfrac2{\sqrt3}$}
\opt{D}{none of these}
\end{qoptions}
\qn{2} Let $A=[a_{ij}]_{3\times3}$ be a $3\times3$ matrix. If every column vector $X$ satisfies $X^TAX=$ Null matrix and $a_{12}=2008$, $a_{13}=1010$, $a_{23}=-2012$. Then the value of $a_{21}+a_{31}+a_{32}$ is (where $X^T$ represents transpose of matrix $X$)
\begin{qoptions}
\opt{A}{$-6$}
\opt{B}{1006}
\opt{C}{$-1006$}
\opt{D}{0}
\end{qoptions}
\qn{3} \textbf{Assertion-(A) :} If the tangent to the curve $y=\dfrac{x-2}{x+1}$ at point $P\left(e,\dfrac{e-2}{e+1}\right)$ makes an angle $\theta$ with positive $x$-axis, then $e^{\cos\theta}<1$.\par
and\par
\textbf{Reason-(R) :} For the curve $xy=-c^2$, $\left.\dfrac{dy}{dx}\right|_{x=t}>0\ \forall\ t\in R$
\begin{qoptions}
\opt{A}{A is True, R is True ; R is a correct explanation for A}
\opt{B}{A is True, R is True ; R is not a correct explanation for A}
\opt{C}{A is True, R is False}
\opt{D}{A is False, R is True}
\end{qoptions}
\qn{4} Consider $f(x)=\dfrac{1-\cos\{x\}}{\left(x^4+ax^3+bx^2+cx+d\right)^2}$ (where $\{.\}$ fractional part of function). If $\displaystyle\lim_{x\to1^+}f(x)$ and $\displaystyle\lim_{x\to2^+}f(x)$ both exist and are non zero then $a+b+c+d$ is equal to
\begin{qoptions}
\opt{A}{2}
\opt{B}{$-1$}
\opt{C}{0}
\opt{D}{5}
\end{qoptions}
\qn{5} Shortest distance between the curve $y=3^x$ \& $y=\log_3x$ is
\begin{qoptions}
\opt{A}{$\left(\dfrac{1+\log(\log3)}{\log3}\right)$}
\opt{B}{$\sqrt2\dfrac{(1+\log(\log3))}{\log3}$}
\opt{C}{$2\sqrt2\left(\dfrac{1+\log(\log3)}{\log3}\right)$}
\opt{D}{None of these}
\end{qoptions}
\qn{6} If the curve $y(x)$ satisfying $x\,dx=\left(\dfrac{x^2}y-y^3\right)dy$ passes through $(0,2)$, then value of $(y(4))^2\left(4-(y(4))^2\right)$ is
\begin{qoptions}
\opt{A}{8}
\opt{B}{16}
\opt{C}{24}
\opt{D}{12}
\end{qoptions}
\qn{7} The value of $\displaystyle\int_{-1/\sqrt3}^{1/\sqrt3}\dfrac{x^4}{1-x^4}\cos^{-1}\left(\dfrac{2x}{1+x^2}\right)dx$ is
\begin{qoptions}
\opt{A}{$\dfrac\pi{12}\left(\pi+3\log\left(2+\sqrt3\right)-4\sqrt3\right)$}
\opt{B}{$\dfrac\pi{12}\left(\pi-3\log\left(2+\sqrt3\right)-4\sqrt3\right)$}
\opt{C}{$\dfrac\pi{12}\left(\pi+3\log\left(2+\sqrt3\right)+4\sqrt3\right)$}
\opt{D}{$\dfrac\pi{12}\left(\pi-3\log\left(2+\sqrt3\right)+4\sqrt3\right)$}
\end{qoptions}
\qn{8} $\displaystyle\lim_{x\to0}\left(2+\log^2_{\sec\left(\frac x2\right)}\cos\left(\dfrac x3\right)\right)^3$ is equal to
\begin{qoptions}
\opt{A}{$\left(\dfrac{146}{81}\right)^3$}
\opt{B}{$\left(\dfrac{70}{27}\right)^3$}
\opt{C}{$\left(\dfrac{178}{81}\right)^3$}
\opt{D}{$\left(\dfrac{34}9\right)^3$}
\end{qoptions}
\qn{9} Match List-I with List-II and select the correct answer using the code given below the list.\par
{\footnotesize\setlength{\tabcolsep}{2pt}\begin{tabular}{|c|p{0.52\linewidth}|c|p{0.08\linewidth}|}
\hline
\multicolumn{2}{|c|}{\textbf{List-I}} & \multicolumn{2}{c|}{\textbf{List-II}}\\ \hline
(A) & If $g(x)=f(-x+xf(x))$ and $f(0)=1$, $f'(0)=2$, then value of $g'(0)$ is equal to & (I) & $\dfrac14$\\ \hline
(B) & If $e^y=e^xx^{-y}$, then $\left(\dfrac{dy}{dx}\right)_{x=e}$ is equal to & (II) & 6\\ \hline
(C) & Number of points of discontinuity of $y=[2x]$ in $[-1,2]$ is equal to (where $[.]$ denote greatest integer function) & (III) & 0\\ \hline
(D) & If $I=\displaystyle\lim_{x\to0}(1+\tan x)^{\frac1{\sin x}}$ then $[I]$ is equal to (where $[.]$ denote greatest integer function) & (IV) & 2\\ \hline
\end{tabular}}\par
\begin{qoptions}
\opt{A}{$A\to III;B\to I;C\to II;D\to IV$}
\opt{B}{$A\to I;B\to II;C\to III;D\to IV$}
\opt{C}{$A\to II;B\to I;C\to IV;D\to III$}
\opt{D}{$A\to II;B\to I;C\to III;D\to IV$}
\end{qoptions}
\qn{10} For $x\in[-2,3]$, the function $f(x)=\left[x^2-2|x|+3\right]$ is [where $[.]$ represents greatest integer function]
\begin{qoptions}
\opt{A}{discontinuous at 4 points}
\opt{B}{discontinuous at 5 points}
\opt{C}{discontinuous at 6 points}
\opt{D}{discontinuous at 7 points}
\end{qoptions}
\qn{11} Consider two functions $f(x)$, $g(x)$ as follows\par
$f(x)=\begin{cases}1-\dfrac{|x|}2 & \text{if}\quad -2\leqslant x\leqslant2\\ 0 & \text{if}\quad x\in(-\infty,-2)\cup(2,\infty)\end{cases}$,\par
$g(x)=f(x+2)+f(x-2)$. The number of points of discontinuity of $g(x)$ is -
\begin{qoptions}
\opt{A}{0}
\opt{B}{1}
\opt{C}{2}
\opt{D}{3}
\end{qoptions}
\qn{12} The focii of an ellipse are located at the points $(2,4)$ and $(2,-2)$. The points $(4,2)$ lies on the ellipse. If $a$ and $b$ represent the lengths of the semi-major and semi-minor axes respectively, then the value of $(ab)^2$ is equal to
\begin{qoptions}
\opt{A}{$68+22\sqrt{10}$}
\opt{B}{$6+22\sqrt{10}$}
\opt{C}{$26+10\sqrt{10}$}
\opt{D}{$6+10\sqrt{10}$}
\end{qoptions}
\qn{13} The differential equation $\left(xy^3(1+\cos x)-y\right)dx+x\,dy=0$ represents the family of curves $\dfrac{x^2}{2y^2}=\dfrac{x^3}{\ell}+x^2\sin x+mx\cos x+n\sin x+C$. Then $\ell+m+n$ is equal to (where $C$ is arbitrary constant)
\begin{qoptions}
\opt{A}{5}
\opt{B}{4}
\opt{C}{1}
\opt{D}{3}
\end{qoptions}
\qn{14} The locus of foot of perpendicular drawn from origin upon any tangent of parabola $y^2=8(x-1)$ is
\begin{qoptions}
\opt{A}{$x^3+xy^2+2y^2+x^2=0$}
\opt{B}{$x^3-xy^2+2y^2-x^2=0$}
\opt{C}{$x^3-xy^2+2y^2+x^2=0$}
\opt{D}{$x^3+xy^2+2y^2-x^2=0$}
\end{qoptions}
\qn{15} Consider the two curves $y=x-\lambda x^2$ and $y=\dfrac{x^2}\lambda$, ($\lambda>0$).\par
\textbf{Statement-1 :} The area bounded between the curves is maximum when $\lambda=1$.\par
and\par
\textbf{Statement-2 :} The area bounded between the curves is $\dfrac{\lambda^2}{(1+\lambda^2)^2}$ square units.
\begin{qoptions}
\opt{A}{Statement-1 is True, Statement-2 is True.}
\opt{B}{Statement-1 is False, Statement-2 is False.}
\opt{C}{Statement-1 is True, Statement-2 is False.}
\opt{D}{Statement-1 is False, Statement-2 is True.}
\end{qoptions}
\qn{16} Two parabolas $x=cy^2+ay+b$ and $x=cy-y^2$ touch each other at the point $(0,1)$, then the value of $2a+b-c$ is-
\begin{qoptions}
\opt{A}{$-1$}
\opt{B}{$-7$}
\opt{C}{$-5$}
\opt{D}{1}
\end{qoptions}
\qn{17} The number of solutions of equation $\left|\ell n|x||-||x|-1|\right|=0$ is
\begin{qoptions}
\opt{A}{2}
\opt{B}{4}
\opt{C}{8}
\opt{D}{10}
\end{qoptions}
\qn{18} The solution set of inequality\par
$\left(\tan^{-1}x\right)\left(\cot^{-1}x\right)-\left(\tan^{-1}x\right)\left(1+\dfrac\pi2\right)-2\cot^{-1}x+2\left(1+\dfrac\pi2\right)>\displaystyle\lim_{x\to-\infty}\left[\sec^{-1}x-\dfrac\pi2\right]$\par
is (where $[x]$ greatest integer function)
\begin{qoptions}
\opt{A}{($\tan1$, $\tan2$)}
\opt{B}{($-\cot1$, $\cot2$)}
\opt{C}{($-\tan1$, $\tan2$)}
\opt{D}{($-\tan1$, $\infty$)}
\end{qoptions}
\qn{19} If $x,y\in R$ and $\sin^{-1}\left(4x^2-4x+2\right)+\cos^{-1}\left(9y^2-6y+2\right)=\dfrac\pi2$ then value of $\dfrac1{x^2}+\dfrac1{y^2}$ is
\begin{qoptions}
\opt{A}{$\dfrac1{13}$}
\opt{B}{25}
\opt{C}{$\dfrac{13}{36}$}
\opt{D}{13}
\end{qoptions}
\qn{20} If $A$ is a matrix $\begin{bmatrix}1&-1&0\\0&2&a\\1&1&2\end{bmatrix}$, then the number of real value(s) of `$a$' for which $A(\mathrm{adj}A)=\mathrm{adj}(\mathrm{adj}\,A)$, is (where $\mathrm{adj}A$ denotes adjoint of matrix $A$)
\begin{qoptions}
\opt{A}{0}
\opt{B}{1}
\opt{C}{2}
\opt{D}{3}
\end{qoptions}
\sect{SECTION-II}
\qn{1} A tall electric pole is to be kept in vertical position by a stretched straight wire from the top pole to the ground. The wire has to clear a wall which is 6 m high and 4 m from the pole. If least length of the wire that can be used between the pole and the ground is $2\left(a^{2/3}+b^{2/3}\right)^{p/q}$ where $a$, $b$, $p$, $q\in N$ and $\dfrac pq$ is in simplest form then the value of $(a+b+p-q)$ is
\qn{2} Number of points of non-differentiability of $f(x)=\left|\left|1-x^3\right|-\left|x^2-1\right|\right|+\left|1-x^2\right|+\left|x^3-1\right|$ is(are)
\qn{3} Let $A$ be the set of all $3\times3$ symmetric matrices whose entries are 1, 1, 1, 0, 0, 0, $-1$, $-1$, $-1$. $B$ is one of the matrix in set $A$ and $X=\begin{bmatrix}x\\y\\z\end{bmatrix}$, $U=\begin{bmatrix}0\\0\\0\end{bmatrix}$. If the number of matrices $B$ such that equation $BX=U$ has infinite solutions is $\lambda$, then $\dfrac\lambda3$ equals
\qn{4} If $\displaystyle\int\dfrac{(x-1)dx}{\left(x+x\sqrt x+\sqrt x\right)\sqrt{\sqrt x(x+1)}}=4\tan^{-1}[g(x)]+C$, where $C$ is an arbitrary constant of integration and $g\!\left(\dfrac12\right)=\sqrt{\dfrac3{\sqrt2}}$, then $(g(1))^2$ is equal to
\qn{5} If $\displaystyle\sum_{m=0}^{2013}\int_0^1x(1-x)^mdx=\dfrac ab$ in lowest form then $b-a$ is
\qn{6} A variable chord of hyperbola $\dfrac{x^2}4-\dfrac{y^2}8=1$ subtends a right angle at center of hyperbola. If this chord touches a fixed circle concentric with hyperbola, then square of radius of this circle is
\qn{7} If $L=\displaystyle\lim_{x\to\infty}x\ell n\left(\dfrac1x\right)-2x\ell n\left(\sin\dfrac1{\sqrt x}\right)$ then $\dfrac1L$ is equal to
\qn{8} The solution of $x^2\dfrac{dy}{dx}-xy=1+\cos\dfrac yx$ is $\tan\left(\dfrac y{px}\right)=c-\dfrac1{qx^2}$ where $c$ is constant, then the value of $(p+q)$ is
\qn{9} Let $P=[p_{ij}]$ be a $3\times3$ non singular matrix and a matrix $Q=[q_{ij}]$ is formed such that $q_{ij}=\displaystyle\sum_{\substack{r=1\\r\neq j}}^{3}p_{ir}$ then value of $\det(P^{-1}Q)$ is
\qn{10} The number of integral values of `$a$' for which $f(x)=\begin{cases}x^3-x^2+10x-5 & ,\ x\leqslant1\\ -2x+\log_2\left(a^2-2\right) & ,\ x>1\end{cases}$ has greatest value at $x=1$, is
\papersec{2}
\partsec{1}{PHYSICS}
\sect{SECTION-I(i)}
\qn{1} A plane electromagnetic wave $\vec E=3000\dfrac Nc\,\hat j\sin\left(\pi x-3\pi\times10^8t\right)$ in $xy$ plane as shown. The induced emf in the loop:
\qfig{p2_phy_s1_q01_1}
\begin{qoptions}
\opt{A}{At $t=0$ is $0$}
\opt{B}{At $t=0.01\ \mu s$ is $12000\ V$}
\opt{C}{At $t=0.003\ \mu s$ is $12000\ V$}
\opt{D}{At $t=\dfrac1{1800}\ \mu s$ is $6000\ V$}
\end{qoptions}
\qfullbegin
\qn{2} Consider a film of thickness $L$ as shown in four different case blow. Notice the observation of the film with light falling normally. Mark the correct statement(s): (Take $L=\dfrac{3\lambda}{2}$, where $\lambda$ is wavelength in $\mu=1.4$)
\qfig{p2_phy_s1_q02_1}
\qfullsplit
\begin{qoptions}
\opt{A}{For (1) and (2), the reflection at film interfaces causes zero phase difference for two reflected rays.}
\opt{B}{For (2) and (3), the reflection at film interfaces causes a phase difference of $\pi$ for two reflected rays.}
\opt{C}{For (1), the film will appear dark, if it is observed through reflected rays from film interfaces.}
\opt{D}{For (3), the film will appear dark, if it is observed through reflected rays from film interfaces}
\end{qoptions}
\qfullend
\qn{3} The following diagram is an arrangement to view the interference pattern produced by a wedge of liquid between a microscope slide and a glass block. The interference pattern obtained is made of equally spaced parallel fringes. The fringe separation (or fringe width) can not be decreased by which of the following actions.
\qfig{p2_phy_s1_q03_1}
\begin{qoptions}
\opt{A}{by increasing the angle of the liquid wedge.}
\opt{B}{by using a liquid of smaller refractive index.}
\opt{C}{by using a thicker glass block.}
\opt{D}{by using a longer liquid wedge of the same angle}
\end{qoptions}
\qn{4} Two narrow cylindrical pipes A and B are filled with hydrogen and helium gases respectively at same temperature. Pipe A is open at both ends and pipe B is open at one end and closed at the other end. If frequency of the second overtone of pipe A is equal to the frequency of the first overtone of the pipe B,
\begin{qoptions}
\opt{A}{Ratio of the velocities of sound waves in the two pipes $\dfrac{v_A}{v_B}=\sqrt2$.}
\opt{B}{Ratio of the velocities of sound waves in the two pipes $\dfrac{v_A}{v_B}=\sqrt{\dfrac{42}{25}}$.}
\opt{C}{Ratio of their lengths $\dfrac{\ell_A}{\ell_B}=\sqrt{\dfrac{168}{25}}$.}
\opt{D}{Ratio of their lengths $\dfrac{\ell_A}{\ell_B}=2\sqrt2$.}
\end{qoptions}
\qfullbegin
\qn{5} A car is chasing truck as shown in figure. If car sounds horn, then choose the CORRECT options.
\qfig{p2_phy_s1_q05_1}
\qfullsplit
\begin{qoptions}
\opt{A}{Apparent frequency of sound for truck driver is more than actual frequency of sound.}
\opt{B}{Apparent wavelength of sound for truck driver is more than actual wavelength of sound.}
\opt{C}{Apparent wavelength of sound for truck driver is less than actual wavelength of sound.}
\opt{D}{Speed of sound observed by truck driver is lesser than actual speed of sound in still air.}
\end{qoptions}
\qfullend
\qn{6} The surface of a spherical star of radius $R$ pulsates. This pulsation can be assumed to be the fundamental mode of vibration with antinode at the surface of the star and node at the centre. If the amplitude of surface is $A$, bulk modulus of the material is $B$ and density is $\rho$, choose the correct statement(s)
\begin{qoptions}
\opt{A}{The frequency of pulsation is $\dfrac1{4R}\sqrt{\dfrac B\rho}$}
\opt{B}{The amplitude at a distance $\dfrac R2$ from the centre is $\dfrac A2$}
\opt{C}{The amplitude at a distance $\dfrac R2$ from the centre is $\dfrac A{\sqrt2}$}
\opt{D}{The maximum velocity of the surface of the star is $\dfrac{\pi A}R\sqrt{\dfrac B\rho}$}
\end{qoptions}
\sect{SECTION-I(ii)}
\qn{7} Match the following graphs (in list-II) between position of object ($u$) and position of image ($v$).\par
{\small\setlength{\tabcolsep}{3pt}\begin{tabular}{|c|p{0.3\linewidth}|c|c|}
\hline
 & \textbf{List--I} & & \textbf{List--II}\\ \hline
(I) & Concave mirror & (P) & \qfigi{p2_phy_s2_q07_1}\\ \hline
(II) & Convex mirror & (Q) & \qfigi{p2_phy_s2_q07_2}\\ \hline
(III) & Convex lens & (R) & \qfigi{p2_phy_s2_q07_3}\\ \hline
(IV) & Concave lens & (S) & \qfigi{p2_phy_s2_q07_4}\\ \hline
 & & (T) & \qfigi{p2_phy_s2_q07_5}\\ \hline
\end{tabular}}\par
\begin{qoptions}
\opt{A}{$\mathrm I\to S;\allowbreak \mathrm{II}\to P;\allowbreak \mathrm{III}\to R;\allowbreak \mathrm{IV}\to Q$}
\opt{B}{$\mathrm I\to R;\allowbreak \mathrm{II}\to Q;\allowbreak \mathrm{III}\to S;\allowbreak \mathrm{IV}\to P$}
\opt{C}{$\mathrm I\to P;\allowbreak \mathrm{II}\to S;\allowbreak \mathrm{III}\to Q;\allowbreak \mathrm{IV}\to R$}
\opt{D}{$\mathrm I\to Q;\allowbreak \mathrm{II}\to R;\allowbreak \mathrm{III}\to P;\allowbreak \mathrm{IV}\to S$}
\end{qoptions}
\qn{8} An object lies on the principal axis of a convex mirror of focal length $f=1.0$ m at a distance $3f$ from it. At time $t=0$, the object starts moving with velocity components $v_0=1$ m/s along the principal axis (towards the mirror) and $v_0=1$ m/s perpendicular to the principal axis. The image moves along the line $ax+by=c$ (where $a$, $b$, $c$ are relatively prime integers and $x$ \& $y$ are in meters.) Match the quantities given in List--I with their magnitudes (in SI units) in List--II. We are only concerned with the time in which the rays from the object to the mirror are paraxial. Origin as well as $x$ \& $y$ axis are considered as per sign convention.
\qfig{p2_phy_s2_q08_1}
{\small\setlength{\tabcolsep}{3pt}\begin{tabular}{|c|p{0.4\linewidth}|c|c|}
\hline
 & \textbf{List--I} & & \textbf{List--II}\\ \hline
(I) & Axial component of velocity of image at time $t=0$ & (P) & $1/3$\\ \hline
(II) & Transverse component of velocity of image at time $t=0$ & (Q) & $1/16$\\ \hline
(III) & The value of $1/a$ & (R) & $1/4$\\ \hline
(IV) & The value of $1/c$ & (S) & $1/6$\\ \hline
 & & (T) & $1/8$\\ \hline
\end{tabular}}\par
\begin{qoptions}
\opt{A}{$\mathrm I\to Q;\allowbreak \mathrm{II}\to R;\allowbreak \mathrm{III}\to R;\allowbreak \mathrm{IV}\to P$}
\opt{B}{$\mathrm I\to P;\allowbreak \mathrm{II}\to R;\allowbreak \mathrm{III}\to S;\allowbreak \mathrm{IV}\to Q$}
\opt{C}{$\mathrm I\to Q;\allowbreak \mathrm{II}\to S;\allowbreak \mathrm{III}\to S;\allowbreak \mathrm{IV}\to P$}
\opt{D}{$\mathrm I\to P;\allowbreak \mathrm{II}\to R;\allowbreak \mathrm{III}\to Q;\allowbreak \mathrm{IV}\to S$}
\end{qoptions}
\qn{9} Match the time period in List-II to the system in list-I.\par
{\small\setlength{\tabcolsep}{2pt}\begin{tabular}{|c|p{0.27\linewidth}|c|c|}
\hline
\multicolumn{2}{|c|}{\textbf{List-I}} & \multicolumn{2}{c|}{\textbf{List-II}}\\ \hline
(I) & A thin hemispherical shell of radius $R$ is supported on the end of thin fixed vertical rod. There is no slipping between rod and shell. \qfigi{p2_phy_s2_q09_1} & (P) & $2\pi\sqrt{\dfrac{26R}{15g}}$\\ \hline
(II) & A thin semicircular ring of radius $R$ is hanging from point $O$ which is accelerating in horizontal direction with some acceleration so that the line joining two ends of the semicircular ring is vertical. \qfigi{p2_phy_s2_q09_2} & (Q) & $2\pi\sqrt{\dfrac{R}{g\sqrt{1+\frac4{\pi^2}}}}$\\ \hline
(III) & Solid hemisphere of radius $R$ is placed on rough sufficiently rough horizontal surface to support rolling. \qfigi{p2_phy_s2_q09_3} & (R) & $2\pi\sqrt{\dfrac{2R}{g\left(1+\frac4{\pi^2}\right)}}$\\ \hline
(IV) & A solid sphere of radius $R$ is placed inside a fixed concave surface of radius $r=6R$. Friction is large between sphere and concave surface to support rolling. \qfigi{p2_phy_s2_q09_4} & (S) & $2\pi\sqrt{\dfrac{4R}{3g}}$\\ \hline
 & & (T) & $2\pi\sqrt{\dfrac{7R}{g}}$\\ \hline
\end{tabular}}\par
Which one of the following options is correct ?
\begin{qoptions}
\opt{A}{$\mathrm I\to T;\allowbreak \mathrm{II}\to Q;\allowbreak \mathrm{III}\to R;\allowbreak \mathrm{IV}\to P$}
\opt{B}{$\mathrm I\to S;\allowbreak \mathrm{II}\to P;\allowbreak \mathrm{III}\to R;\allowbreak \mathrm{IV}\to Q$}
\opt{C}{$\mathrm I\to S;\allowbreak \mathrm{II}\to R;\allowbreak \mathrm{III}\to T;\allowbreak \mathrm{IV}\to P$}
\opt{D}{$\mathrm I\to S;\allowbreak \mathrm{II}\to R;\allowbreak \mathrm{III}\to P;\allowbreak \mathrm{IV}\to T$}
\end{qoptions}
\qn{10} In resonance column experiment, successive resonance occurs for pipe of length $L_1=40$ cm, $L_2=67.4$ cm. [$c=320$ m/s]. Match approximately List--I with List--II.\par
{\small\setlength{\tabcolsep}{3pt}\begin{tabular}{|c|p{0.4\linewidth}|c|c|}
\hline
 & \textbf{List--I} & & \textbf{List--II}\\ \hline
(I) & Ratio of harmonics that $L_1$ \& $L_2$ corresponds to & (P) & $12.6$\\ \hline
(II) & end correction (in millimeter) & (Q) & $1.2$\\ \hline
(III) & minimum tube length for which resonance can occur (in cm) & (R) & $0.6$\\ \hline
(IV) & Ratio of pressure amplitude at $x=\lambda/3$ to $x=\lambda/8$ from the closed end when tube length is $L_1$. & (S) & $0.8$\\ \hline
 & & (T) & $11$\\ \hline
\end{tabular}}\par
\begin{qoptions}
\opt{A}{$\mathrm I\to S;\allowbreak \mathrm{II}\to P;\allowbreak \mathrm{III}\to R;\allowbreak \mathrm{IV}\to Q$}
\opt{B}{$\mathrm I\to R;\allowbreak \mathrm{II}\to T;\allowbreak \mathrm{III}\to P;\allowbreak \mathrm{IV}\to Q$}
\opt{C}{$\mathrm I\to P;\allowbreak \mathrm{II}\to S;\allowbreak \mathrm{III}\to Q;\allowbreak \mathrm{IV}\to R$}
\opt{D}{$\mathrm I\to P;\allowbreak \mathrm{II}\to Q;\allowbreak \mathrm{III}\to R;\allowbreak \mathrm{IV}\to S$}
\end{qoptions}
\sect{SECTION-II}
\qn{1} If you place the flat side of a hemispherical glass object (lens) against a paper, it is possible to see a magnified image of the paper's surface around the center of the lens. How many times is the image bigger than the object if you are observing at a distance that is much bigger than the dimensions of the lens? The refractive index of glass is $n=1.5$.
\qn{2} A straight conducting wire $MN$ of length 4m is rotated in a uniform magnetic field $\vec B=0.1$ Tesla $(\hat k)$ about point $O$ with angular speed $\omega=10$ rad/s as shown in the figure. The axis of rotation is parallel to $\vec B$. Find the potential difference (in volt) between point $M$ and point $N$.
\qfig{p2_phy_s3_q02_1}
\qn{3} The reading of the ammeter in the circuit in figure is
\qfig{p2_phy_s3_q03_1}
\qn{4} A steady force of 120 N is required to push a boat of mass 700 kg through water at a constant speed of 1 m/s. If the boat is fastened by a spring and held at 2m from the equilibrium position by a force of 450 N, find the angular frequency (in rad/s) of damped SHM.
\qfig{p2_phy_s3_q04_1}
\qn{5} We use a simple microscope to magnify an object. The microscope has a numerical aperture of $\sin\alpha=0.24$. The object is so small that the resolving power of the microscope is totally employed. If the diameter of the eye's pupil is $d=4.0$ mm, find the minimum magnifying power of the microscope ?
\qn{6} In the given network all the capacitors have same capacitance 1micro farad. Find equivalent capacitance between A and B in micro farad.
\qfig{p2_phy_s3_q06_1}
\qn{7} Figure shows the arrangement of a set of saddle coils. If the magnitude of magnetic field at point P is $\dfrac{\alpha\mu_0i}{\pi a}$, then value of $\alpha$ is
\qfig{p2_phy_s3_q07_1}
\qn{8} A string of mass per unit length $\mu$ is clamped at both ends such that one end of the string is at $x=0$ and the other end at $x=L$. When string vibrates in fundamental mode, amplitude of the midpoint of string is $a$ and tension in string is $F$. Find the total oscillation energy (in J) stored in the string. (Use $L=1$m, $F=10$ N, $a=\dfrac1\pi$m)
\qfig{p2_phy_s3_q08_1}
\partsec{2}{CHEMISTRY}
\sect{SECTION-I(i)}
\qn{1} A hydrated green salt (A) loses approx 45\% of its weight on heating and rendered anhydrous. (A) on strong heating yields a dirty brown solid (B) with evolution of two gases (C) and (D). Gas (C) turned orange colour of dichromate solution into green. Dirty brown solid (B) dissolved in boiling HCl solution to produce a yellow solution of (E). Solution (E) was separated into two parts and one part was treated with $\mathrm{NH_4SCN}$ solution giving a blood-red solution. To the other part of solution, $\mathrm{H_2S}$ was passed , where a white turbidity was observed initially and the solution became light green. An aqueous solution of salt (A) gave white precipitate when treated with aqueous barium chloride solution.\par
Consider above given paragraph and select correct option :-
\begin{qoptions}
\opt{A}{Heating of anhydrous salt (A) is an intramolecular redox reaction.}
\opt{B}{C = $\mathrm{SO_2}$}
\opt{C}{D = $\mathrm{SO_3}$}
\opt{D}{Reagent (E) gives blue colour solution with $\mathrm{K_3[Fe(CN)_6]}$}
\end{qoptions}
\qn{2} In which of the following reaction(s) $\mathrm{H_2O_2}$ act as reducing agent?
\begin{qoptions}
\opt{A}{$\mathrm{PbS+4H_2O_2\longrightarrow PbSO_4+4H_2O}$}
\opt{B}{{\footnotesize$\mathrm{2KMnO_4+3H_2SO_4+5H_2O_2}\longrightarrow$ $\mathrm{2MnSO_4+K_2SO_4+}$ $\mathrm{8H_2O+5O_2}$}}
\opt{C}{$\mathrm{HOCl+H_2O_2\longrightarrow HCl+H_2O+O_2}$}
\opt{D}{$\mathrm{I_2+H_2O_2+2KOH\longrightarrow 2KI+2H_2O+O_2}$}
\end{qoptions}
\qn{3} Select the correct statement(s) :
\begin{qoptions}
\opt{A}{In irreversible adiabatic compression of ideal gas entropy of system increases.}
\opt{B}{Diffusion of gas from high pressure to low pressure is spontaneous phenomenon.}
\opt{C}{In a reversible chemical process, entropy change of universe is positive.}
\opt{D}{A chemical reaction is always spontaneous if entropy of universe increases due to it.}
\end{qoptions}
\qn{4} pH of 0.1 M weak monoprotic acid (HA) is 1.7 at $25^\circ$C. ($\log2=0.30$, $\log10=1$)
\begin{qoptions}
\opt{A}{Acid dissociation constant of HA is $5\times10^{-3}$.}
\opt{B}{Degree of dissociation of 0.1 M HA solution is 20\%.}
\opt{C}{pOH of 0.1 M HA solution is 11.3.}
\opt{D}{Concentration of $\mathrm{H^+}$ ion in solution decreases on dilution of solution.}
\end{qoptions}
\qfullbegin
\qn{5}
\qfig{p2_chem_s1_q05_1}
\qfullsplit
If T can evolve effervescence of $\mathrm{CO_2}$ with $\mathrm{NaHCO_3}$ solution then correct statement(s) is/are :
\begin{qoptions}
\opt{A}{Heating of `T' will release water}
\opt{B}{T is most acidic among all isomeric benzenoid dicarboxylic acid}
\opt{C}{Q \& S can be distinguished by mustered oil test}
\opt{D}{T can give benzene as major product with sodalime / $\Delta$}
\end{qoptions}
\qfullend
\qn{6} Among P, Q, R and S, the aromatic compound(s) is / are :
\qfig{p2_chem_s1_q06_1}
\qfig{p2_chem_s1_q06_2}
\begin{qoptions}
\opt{A}{P}
\opt{B}{Q}
\opt{C}{R}
\opt{D}{S}
\end{qoptions}
\sect{SECTION-I(ii)}
\qn{7} Match each coordination compound in List-I with an appropriate pair of characteristics from List-II and select the correct answer using the code given below the lists. \{en = $\mathrm{H_2NCH_2CH_2NH_2}$ ' atomic numbers ; Ti = 22 ; Cr = 24 ; Co = 27 ; Pt = 78\}\par
{\footnotesize\setlength{\tabcolsep}{2pt}\begin{tabular}{|c|p{0.34\linewidth}|c|p{0.34\linewidth}|}
\hline
\multicolumn{2}{|c|}{\textbf{List-I}} & \multicolumn{2}{c|}{\textbf{List-II}}\\ \hline
(P) & Tetraammine\-dichlorido\-chromium(III) chloride & (1) & Para\-magnetic and exhibits ionisation iso\-merism\\ \hline
(Q) & Pentaaqua\-chlorido\-titanium(III) nitrate & (2) & Di\-magnetic and exhibits cis-trans isomerism\\ \hline
(R) & Ammine\-chlorido(ethylene\-diammine)\-platinum(II) nitrate & (3) & Para\-magnetic and exhibits cis-trans isomerism\\ \hline
(S) & Tetraammine\-dinitrato\-cobalt(III) nitrate & (4) & Dia\-magnetic and exhibit ionisation iso\-merism\\ \hline
 & & (5) & Low spin complex\\ \hline
\end{tabular}}\par
The correct option is :
\begin{qoptions}
\opt{A}{$P\to3;\allowbreak Q\to2;\allowbreak R\to3,5;\allowbreak S\to3,5$}
\opt{B}{$P\to3;\allowbreak Q\to1;\allowbreak R\to4,5;\allowbreak S\to2,5$}
\opt{C}{$P\to3,5;\allowbreak Q\to1;\allowbreak R\to3;\allowbreak S\to4$}
\opt{D}{$P\to1;\allowbreak Q\to3,5;\allowbreak R\to4;\allowbreak S\to2$}
\end{qoptions}
\qn{8} Match the correct characteristics for products produced by reactions mentioned in List-I.\par
{\small\setlength{\tabcolsep}{3pt}\begin{tabular}{|c|p{0.3\linewidth}|c|p{0.32\linewidth}|}
\hline
\multicolumn{2}{|c|}{\textbf{List-I} Reactions} & \multicolumn{2}{c|}{\textbf{List-II} Characteristics}\\ \hline
(P) & $\mathrm{Pb^{2+}+H_2S}$ & (1) & Disproportionation reaction\\ \hline
(Q) & $\mathrm{Au+}$ Aqua regia & (2) & Intermolecular redox reaction\\ \hline
(R) & $\mathrm{Cu+}$ conc. $\mathrm{H_2SO_4}$ & (3) & Precipitate formation\\ \hline
(S) & $\mathrm{AgCl+}$ Ammonia solution & (4) & Gaseous product formed\\ \hline
 & & (5) & $+1$ oxidation state of metal ion\\ \hline
\end{tabular}}\par
\begin{qoptions}
\opt{A}{$P\to1,5;\allowbreak Q\to2;\allowbreak R\to2,4;\allowbreak S\to1,4,5$}
\opt{B}{$P\to1,4;\allowbreak Q\to2,4;\allowbreak R\to4;\allowbreak S\to1,4,5$}
\opt{C}{$P\to3;\allowbreak Q\to4;\allowbreak R\to2,4;\allowbreak S\to1,3,5$}
\opt{D}{$P\to3;\allowbreak Q\to2,4;\allowbreak R\to2,4;\allowbreak S\to5$}
\end{qoptions}
\qn{9} Match the list.\par
{\footnotesize\setlength{\tabcolsep}{2pt}\begin{tabular}{|c|p{0.38\linewidth}|c|p{0.3\linewidth}|}
\hline
 & \textbf{List-I} & & \textbf{List-II}\\ \hline
(P) & $\mathrm{N_2O_{4(g)}\rightleftharpoons2NO_{2(g)}}$ on increasing volume at constant temperature & (1) & Forward shift.\\ \hline
(Q) & $\mathrm{H_{2(g)}+I_{2(g)}\rightleftharpoons2HI_{(g)}}$ on addition of inert gas at constant P \& T. & (2) & Concentration of reactant decreases.\\ \hline
(R) & $\mathrm{2NO_{(g)}+Br_{2(g)}}$ $\rightleftharpoons$ $\mathrm{2NOBr_{(g)}}$ on decreasing pressure at constant temperature. & (3) & Concentration of product decreases.\\ \hline
(S) & $\mathrm{2ICl_{3(g)}}$ $\rightleftharpoons$ $\mathrm{I_{2(g)}+3Cl_{2(g)}}$ $[\Delta H>0]$ on increasing pressure at constant temperature. & (4) & Equilibrium is not disturbed.\\ \hline
 & & (5) & $K_P>K_C\ \left(T>\dfrac1R\right)$\\ \hline
\end{tabular}}\par
\begin{qoptions}
\opt{A}{$P\to1,2,5;\allowbreak Q\to2,3,4;\allowbreak R\to2,3;\allowbreak S\to3,5$}
\opt{B}{$P\to1,2,5;\allowbreak Q\to1,5;\allowbreak R\to2,3;\allowbreak S\to3,5$}
\opt{C}{$P\to3;\allowbreak Q\to1,4;\allowbreak R\to1;\allowbreak S\to2$}
\opt{D}{$P\to1,2,3,5;\allowbreak Q\to2,3,4;\allowbreak R\to2,3;\allowbreak S\to5$}
\end{qoptions}
\qn{10}
{\footnotesize\setlength{\tabcolsep}{2pt}\begin{tabular}{|c|p{0.25\linewidth}|c|p{0.28\linewidth}|}
\hline
\multicolumn{2}{|c|}{\textbf{List - I} (Compound)} & \multicolumn{2}{c|}{\textbf{List - II} (Can be distinguish by)}\\ \hline
(P) & $\mathrm{CH_3{-}CHO}$ \& $\mathrm{Ph{-}CHO}$ & (1) & $\mathrm{Br_2/H_2O}$ test\\ \hline
(Q) & $\mathrm{CH_3{-}CH_3}$ \& $\mathrm{HC{\equiv}CH}$ & (2) & Ammonical $\mathrm{AgNO_3}$\\ \hline
(R) & $\mathrm{CH_3-C{\equiv}CH}$ \& \qfigi{p2_chem_s2_q09_1} & (3) & Iodoform test\\ \hline
(S) & $\mathrm{CH_3{-}CHO}$ \& \qfigi{p2_chem_s2_q09_2} & (4) & $\mathrm{NaHSO_3}$ test\\ \hline
 & & (5) & Fehling test\\ \hline
\end{tabular}}\par
Choose correct option?
\begin{qoptions}
\opt{A}{$P\to2,5;\allowbreak Q\to3,5;\allowbreak R\to1,5;\allowbreak S\to2,5$}
\opt{B}{$P\to3,5;\allowbreak Q\to1,2;\allowbreak R\to1,2,3,4;\allowbreak S\to2,5$}
\opt{C}{$P\to3,5;\allowbreak Q\to1,2;\allowbreak R\to5;\allowbreak S\to3,5$}
\opt{D}{$P\to2,5;\allowbreak Q\to3,5;\allowbreak R\to1,2,3,4;\allowbreak S\to3,5$}
\end{qoptions}
\sect{SECTION-II}
\qn{1} The compound KCN is traditionally used to extract gold from ore. Gold dissolves in cyanide solution in the presence of air to form $\mathrm{Au(CN)_2^-}$.\par
How many grams of KCN are needed to extract 20 gram of gold from ore?\par
(Atomic weight of Au = 197 , K = 39 , C = 12 , N = 14)
\qn{2} Consider the following compounds :\par
$\mathrm{H_2S_2O_3}$ , $\mathrm{H_2S_2O_4}$ , $\mathrm{H_2S_2O_5}$ , $\mathrm{H_2S_2O_6}$ , $\mathrm{H_2S_2O_7}$ , $\mathrm{H_2S_2O_8}$\par
Calculate value of $(P\times Q\times R)$\par
where $P$ = Number of oxyacids in which all sulphur is in $+6$ oxidation state\par
$Q$ = Number of oxyacids in which all Sulpur atom bonded by 6 bonds\par
$R$ = Number of oxyacids in which at least one S-O-S linkage present.
\qn{3} The progress of reaction, $\mathrm{A(g)\rightleftharpoons xB(g)+yC(g)}$ with time is presented in figure. The absolute value of change in standard gibbs energy of reaction is $|\Delta_rG^\circ|$, (in cal/mole) at 300 K. Find $\dfrac{|\Delta_rG^\circ|}{140}$.\par
Given : $R=2$ cal/mole-K, $\ln2=0.7$
\qfig{p2_chem_s3_q03_1}
\qn{4} The entropy versus temperature plot for phases $\alpha$ and $\beta$ at 1 bar pressure is given. $S_T$ and $S_0$ are entropies of the phases at temperatures $T$ and 0 K, respectively.
\qfig{p2_chem_s3_q04_1}
The transition temperature for $\alpha$ to $\beta$ phase change is 600 K and $C_{p,\beta}-C_{p,\alpha}=1\,\mathrm{J\,mol^{-1}\,K^{-1}}$. Assume $(C_{p,\beta}-C_{p,\alpha})$ is independent of temperature in the range of 200 to 700 K. $C_{p,\alpha}$ and $C_{p,\beta}$ are heat capacities of $\alpha$ and $\beta$ phases, respectively.\par
The value of $\Delta G$ for the transition $\alpha\longrightarrow\beta$ at 300 K in joule is ..........\par
[Use : $\ln2=0.69$, Given : $S_\beta-S_\alpha=0$ at 0 K]
\qn{5} The vapour pressure of two miscible liquids (A) and (B) are 300 and 500 mm of Hg respectively. In a flask 10 mole of (A) is mixed with 12 mole of (B). However, as soon as (B) is added, (A) starts polymerising into a completely insoluble solid. The polymerisation follows first-order kinetics. After 100 minute, 0.525 mole of a solute is dissolved which arrests the polymerisation completely. The final vapour pressure of the solution is 400 mm of Hg. Estimate the rate constant of the polymerisation reaction. Assume negligible volume change on mixing and polymerisation and ideal behaviour for the final solution. ($\ln9.9=2.29$)\par
If your answer is $x\times10^{-y}$ (in $\mathrm{min^{-1}}$) in scientific notation then `y' is.
\qfullbegin
\qn{6} Total number of amine among following which can be prepared in good yield by Gabriel phthalimide synthesis method = $X$.\par
Total number of amines giving isocyanide test = $Y$.
\qfig{p2_chem_s3_q06_1}
\qfullsplit
Then what will value of $Y-X$ :
\qfullend
\qn{7} Find total number of reaction in which number of carbon atom decreases in final major product?
\qfig{p2_chem_s3_q07_1}
\qfig{p2_chem_s3_q07_2}
\qfig{p2_chem_s3_q07_3}
\qn{8} Consider following reaction\par
(i) $\mathrm{CH_3-C{\equiv}C-H}\xrightarrow[\text{Fe tube}/\Delta]{\text{Red hot}}(\mathrm A)\xrightarrow{\mathrm{KMnO_4/H^\oplus}}(\mathrm B)$ \quad ((A) Major)\par
\qfigi{p2_chem_s3_q08_1}$\xrightarrow{\mathrm{H_2/Pd-BaSO_4}}(\mathrm P)$\par
$(\mathrm P)\xrightarrow[\mathrm{(ii)\ H^\oplus/H_2O/\Delta}]{\mathrm{(i)(CH_3CO)_2O/CH_3COONa,\ \Delta}}\mathrm Q$ \quad (Major)\par
If degree of unsaturation of $\mathrm B$ is $x$ and geometrical isomer of $\mathrm Q$ is $y$ then find $x+y$?
\partsec{3}{MATHEMATICS}
\sect{SECTION-I(i)}
\qn{1} If point $(4,4)$ lies on parabola whose focus lies on $x$-axis and directrix is $\ell x+my=1$ such that $(\ell,m)$ lies on curve $x^2+y^2-32xy+8x+8y-1=0$ then :
\begin{qoptions}
\opt{A}{co-ordinates of focus may be $(5,0)$}
\opt{B}{co-ordinates of focus may be $(3,0)$}
\opt{C}{If tangent at $(4,4)$ passes through origin, value of $m$ may be $\dfrac13$}
\opt{D}{If tangent at $(4,4)$ passes through origin, equation of axis may be $x-4y-3=0$}
\end{qoptions}
\qn{2} Let $M$ be a $3\times3$ matrix satisfying $M^3=0$, then which of the following is/are true -
\begin{qoptions}
\opt{A}{$|M^2-M+I|\neq0$}
\opt{B}{$|M^2-M+I|=0$}
\opt{C}{$\left|\dfrac{M^2}2+M+I\right|\neq0$}
\opt{D}{$\left|\dfrac{M^2}2+M+I\right|=0$}
\end{qoptions}
\qn{3} Let $f(x)=\begin{cases}4x^2+1 & x<\frac12\\ 2(\log_ex+1) & \frac12\le x<1\\ x^2+1 & x\ge1\end{cases}$, then which of the following is/are true ?
\begin{qoptions}
\opt{A}{$f(x)$ has two points of extremum}
\opt{B}{$f(x)$ is continuous $\forall\ x\in R$, except at one point}
\opt{C}{$f(x)$ is differentiable $\forall\ x\in R$ except at one point}
\opt{D}{$f(x)$ has exactly one point of minima}
\end{qoptions}
\qn{4} Let $f(x)$ be a non-constant, twice differentiable function defined on $R$ such that $y=f(x)$ is symmetric about line $x=1$ and $f(-1)=f'\left(\dfrac14\right)=f'\left(\dfrac12\right)=0$, then which of following is/are correct statement(s) ?
\begin{qoptions}
\opt{A}{$f''(x)=0$ has atleast four roots in the interval $(0,2)$}
\opt{B}{$\displaystyle\int_{-2}^{2}(x^5+x^3)f(1+x)\,dx<\int_0^2\dfrac1{1+2^{f(x)}}dx$}
\opt{C}{there exist at least one $C\in\left(\dfrac32,\dfrac74\right)$ such that $f'(C)+Cf''(C)=0$}
\opt{D}{for some $C\in\left(1,\dfrac32\right)$, $f'(C)=Cf''(C)$}
\end{qoptions}
\qn{5} Let $A$ is a non-singular matrix of order $n\times n$ such that $3ABA^{-1}+A=2A^{-1}BA$, then which is/are correct (where $|A|=\det(A)$) -
\begin{qoptions}
\opt{A}{$A$ and $B$ both are identity matrix}
\opt{B}{$|A+B|=0$}
\opt{C}{$|ABA^{-1}-A^{-1}BA|=0$}
\opt{D}{$A+B$ is not a singular matrix}
\end{qoptions}
\qn{6} Consider an ellipse with major axis and minor axis having length 4 and $2\sqrt3$ respectively. Let $P$ be variable point on it and $S$ and $S'$ be the focii and $C$ be the centre of ellipse then
\begin{qoptions}
\opt{A}{maximum value of $PS.PS'$ is 4}
\opt{B}{minimum value of $PS^2+PS'^2$ is 8}
\opt{C}{minimum length of portion of tangent at $P$ intercepted between the principal axes is $\sqrt{7-4\sqrt3}$}
\opt{D}{If the foot of perpendicular from centre upon the normal at point $P$ is the point $Q$ then the maximum value of $PQ$ is $2-\sqrt3$}
\end{qoptions}
\sect{SECTION-I(ii)}
\qn{7}
{\small\setlength{\tabcolsep}{2pt}\begin{tabular}{|c|p{0.58\linewidth}|c|c|}
\hline
\multicolumn{2}{|c|}{\textbf{List-I}} & \multicolumn{2}{c|}{\textbf{List-II}}\\ \hline
(I) & Value of $\lim\limits_{x\to0}\log_{\frac{x^2+1}{x+1}}\left(\sqrt{x^2+1}-x\right)$ is less than or equal to & (P) & $-1$\\ \hline
(II) & The value of $\lim\limits_{x\to\infty}\dfrac{2\left(3^{\sin x}+2x+1\right)}{\sin x+\sqrt{x^2+1}}$ is greater than or equal to & (Q) & $4$\\ \hline
(III) & Let $f:R\to\left[0,\dfrac\pi2\right)$ be defined by $f(x)=\tan^{-1}\left(3x^2+6x+a\right)$. If $f(x)$ is an onto function then the value of $a$ can be & (R) & $3$\\ \hline
(IV) & Let $x^5+10x^3y^2+15x^2y^3+3x^4y+y^5=0$ and $\dfrac{dy}{dx}=\dfrac{\lambda y}x$, then $\lambda$ is equal to & (S) & $5$\\ \hline
 & & (T) & $1$\\ \hline
\end{tabular}}\par
\begin{qoptions}
\opt{A}{$\mathrm I\to Q,R,S,T;\allowbreak \mathrm{II}\to P,Q,R,T;\allowbreak \mathrm{III}\to R;\allowbreak \mathrm{IV}\to T$}
\opt{B}{$\mathrm I\to P,Q,R,T;\allowbreak \mathrm{II}\to Q,R,S,T;\allowbreak \mathrm{III}\to R;\allowbreak \mathrm{IV}\to T$}
\opt{C}{$\mathrm I\to Q,R,S,T;\allowbreak \mathrm{II}\to R;\allowbreak \mathrm{III}\to P,Q,R,T;\allowbreak \mathrm{IV}\to T$}
\opt{D}{$\mathrm I\to Q,R,S,T;\allowbreak \mathrm{II}\to P,Q,R,T;\allowbreak \mathrm{III}\to T;\allowbreak \mathrm{IV}\to R$}
\end{qoptions}
\qn{8}
{\footnotesize\setlength{\tabcolsep}{2pt}\begin{tabular}{|c|p{0.34\linewidth}|c|p{0.27\linewidth}|}
\hline
\multicolumn{2}{|c|}{\textbf{List-I}} & \multicolumn{2}{c|}{\textbf{List-II}}\\ \hline
(I) & $f(x)=\left[\dfrac{4x}\pi\right]$ & (P) & range contains atleast one integer but not more than seven\\ \hline
(II) & $f(x)=\cos^{-1}$ $\left(\mathrm{sgn}\left(\cos\dfrac{2x-1}2\pi\right)\right)$ & (Q) & is many one in $[-2,2]$\\ \hline
(III) & $f(x)=\max(\{x+1\},\{5-x\})$ & (R) & is even function in $[-2,2]$\\ \hline
(IV) & $f(x)=\sqrt{x^2}+[x]^2$ & (S) & is neither odd nor even function in $[-2,2]$\\ \hline
 & & (T) & is bounded function\\ \hline
\end{tabular}}\par
[\textbf{Note :} where $[\cdot]$ denotes greatest integer function, $\{\cdot\}$ denotes fractional part function and $\mathrm{sgn}\,(\cdot)$ denotes signum function]
\begin{qoptions}
\opt{A}{$\mathrm I\to Q,S;\allowbreak \mathrm{II}\to P,Q,S,T;\allowbreak \mathrm{III}\to Q,S;\allowbreak \mathrm{IV}\to Q,R,T$}
\opt{B}{$\mathrm I\to Q,S;\allowbreak \mathrm{II}\to P,Q,S,T;\allowbreak \mathrm{III}\to Q,R,T;\allowbreak \mathrm{IV}\to Q,S$}
\opt{C}{$\mathrm I\to P,Q,S,T;\allowbreak \mathrm{II}\to Q,S;\allowbreak \mathrm{III}\to Q,R,T;\allowbreak \mathrm{IV}\to Q,S$}
\opt{D}{$\mathrm I\to Q,S;\allowbreak \mathrm{II}\to Q,R,T;\allowbreak \mathrm{III}\to P,Q,S,T;\allowbreak \mathrm{IV}\to Q,S$}
\end{qoptions}
\qn{9} Match List-I with List-II and select the correct answer using the code given below the list.\par
{\small\setlength{\tabcolsep}{2pt}\begin{tabular}{|c|p{0.58\linewidth}|c|c|}
\hline
\multicolumn{2}{|c|}{\textbf{List-I}} & \multicolumn{2}{c|}{\textbf{List-II}}\\ \hline
(I) & The value of definite integral $\displaystyle\int_0^{\pi/2}\dfrac{(\sin x)^{\cos x}}{(\sin x)^{\cos x}+(\cos x)^{\sin x}}dx$ is less than & (P) & $1$\\ \hline
(II) & $L=\dfrac{\displaystyle\int_0^{51}[x+25]\,dx}{\displaystyle\int_0^{51}\{x+25\}\,dx}$ then value of $\dfrac L{10}$ is greater than (where $[.]$ \& $\{x\}$ greatest integer function and fractional part function respetively) & (Q) & $2$\\ \hline
(III) & If $f(x)=\left|\left(x^3+bx^2+cx+d\right)\right.$ $\left.(x-1)(x-2)(x-3)^3(x-4)^4\right|$ is derivable for all $x\in R$ and $f'(3)+f'''(4)=0$ then value of $|b+c+d|$ is less than & (R) & $3$\\ \hline
(IV) & If $\displaystyle\lim_{t\to0}\left[\min_{x\in R}\left(x^2+2x-2\right)\dfrac{\sin^{-1}(t)}t\right]$ where $[.]$ greatest integer function then value of $7+L$ is greater than & (S) & $4$\\ \hline
 & & (T) & $5$\\ \hline
\end{tabular}}\par
\begin{qoptions}
\opt{A}{$\mathrm I\to P,Q,R,S,T;\allowbreak \mathrm{II}\to P,Q,R,S;\allowbreak \mathrm{III}\to Q,R,S,T;\allowbreak \mathrm{IV}\to P,Q$}
\opt{B}{$\mathrm I\to P,Q,R,T;\allowbreak \mathrm{II}\to P,Q,R,S;\allowbreak \mathrm{III}\to Q,R,T;\allowbreak \mathrm{IV}\to P,Q$}
\opt{C}{$\mathrm I\to P,R,S,T;\allowbreak \mathrm{II}\to P,Q,R,T;\allowbreak \mathrm{III}\to P,R,T;\allowbreak \mathrm{IV}\to P,R$}
\opt{D}{$\mathrm I\to P,Q,R,S,T;\allowbreak \mathrm{II}\to P,R,S;\allowbreak \mathrm{III}\to R,S,T;\allowbreak \mathrm{IV}\to P,Q$}
\end{qoptions}
\qn{10}
{\small\setlength{\tabcolsep}{2pt}\begin{tabular}{|c|p{0.58\linewidth}|c|c|}
\hline
\multicolumn{2}{|c|}{\textbf{List-I}} & \multicolumn{2}{c|}{\textbf{List-II}}\\ \hline
(I) & The area bounded by the curves $y=\tan^2x$, $y=\sec^2x$, $x=-\dfrac\pi4$ and $x=\dfrac\pi4$ is given by & (P) & $\dfrac{16}3$\\ \hline
(II) & The area bounded by $y=x^2-4$ with the co-ordinate axes is & (Q) & $\dfrac\pi2$\\ \hline
(III) & $\displaystyle\int_0^{\pi/2}\dfrac{\sin6x}{\sin x}dx$ is equal to & (R) & $\dfrac{32}3$\\ \hline
(IV) & The slope of line which cuts-off equal intercept on co-ordinate axes can be & (S) & $\dfrac{26}{15}$\\ \hline
 & & (T) & $1$\\ \hline
\end{tabular}}\par
\begin{qoptions}
\opt{A}{$\mathrm I\to P,Q,R,S,T;\allowbreak \mathrm{II}\to S;\allowbreak \mathrm{III}\to P,R;\allowbreak \mathrm{IV}\to Q$}
\opt{B}{$\mathrm I\to Q;\allowbreak \mathrm{II}\to P,R;\allowbreak \mathrm{III}\to P,Q,R,S,T;\allowbreak \mathrm{IV}\to S$}
\opt{C}{$\mathrm I\to P,R;\allowbreak \mathrm{II}\to Q;\allowbreak \mathrm{III}\to S;\allowbreak \mathrm{IV}\to P,Q,R,S,T$}
\opt{D}{$\mathrm I\to Q;\allowbreak \mathrm{II}\to P,R;\allowbreak \mathrm{III}\to S;\allowbreak \mathrm{IV}\to P,Q,R,S,T$}
\end{qoptions}
\sect{SECTION-II}
\qn{1} Let $P(x)$ be a polynomial of degree 4 with $P(1)=1$, $P'(1)=0$, $P''(1)=2$, $P'''(1)=-12$ and $P''''(1)=24$ then value of $\dfrac{P(2)}{P(3)}$ is equal to
\qn{2} $\displaystyle\lim_{x\to0^+}(x\cos x)^x+(x\sin x)^{1/x}$ is
\qn{3} Let $f(x)=\begin{cases}x^2+1, & x<0\\ (x-1)^3+2a, & x\ge0\end{cases}$ and $g(x)=\begin{cases}x+2b, & x<0\\ |x-2|, & x\ge0\end{cases}$ where $a\in R$ and $b\in R^-$. If $fog(x)$ is continuous for all real $x$ and differentiable at $x=0$ then the value of $[a-b]$ is (where $[\cdot]$ is greatest integer function).
\qn{4} If $\dfrac{\displaystyle\int_0^\pi x^3\ln(\sin x)dx}{\displaystyle\int_0^\pi x^2\ln\left(\sqrt2\sin x\right)dx}=k\pi$, then $k$ is equal to
\qn{5} If $\displaystyle\int\dfrac{dx}{(\sec x+\cos\mathrm{ec}\,x+\tan x+\cot x)^2}dx=\dfrac x\alpha+\dfrac{\cos\left(x+\frac\pi4\right)}\beta+\dfrac{\cos2x}\gamma+\delta$, then $|\alpha+\beta+\gamma|$ is
\qn{6} Let $\lambda=\log_e\left(\displaystyle\int_0^1e^{x^2+x}(2x^2+x+1)\,dx\right)$, the value of $\dfrac\lambda4$ equals
\qn{7} Let $f(x)=|||x|-1|-1|$ and $g(x)=\{\max f(t);\ x\le t\le x+1,\ x\in R\}$. The number of points of non-differentiability of $g(x)$ is
\qn{8} Let $x^2+3y^2=3$ and $m$, $M$ denote smallest and greatest value of expression $(x+y+3)^2-9$ respectively, then the value of $M-m$ is
\end{multicols}
\end{document}
